\documentclass[12pt, oneside]{article}

\usepackage{xparse}

\usepackage{geometry}
\geometry{letterpaper}
\usepackage[parfill]{parskip}
\usepackage{float}

% Math packages
\usepackage{amssymb}
\usepackage{amsmath}
\usepackage{amsthm}
\usepackage{mathalpha}

% Tables
\usepackage{array}
\usepackage{siunitx}
\usepackage{booktabs}
\usepackage{afterpage}

% Hyperlinks
\usepackage{hyperref}

% References
\numberwithin{equation}{section}
\usepackage[numbers]{natbib}

% Theorem environments
\theoremstyle{definition}
\newtheorem{definition}{Definition}[section]
\newtheorem{proposition}{Proposition}[section]
\newtheorem{axiom}{Axiom}[section]

\theoremstyle{plain}
\newtheorem{theorem}{Theorem}[section]
\newtheorem{lemma}[theorem]{Lemma}
\newtheorem{corollary}[theorem]{Corollary}
\newtheorem{conjecture}[theorem]{Conjecture}
\newtheorem{example}[theorem]{Example}

\theoremstyle{remark}
\newtheorem*{remark}{Remark}

\title{A Symbolic--Affine Framework for Density-One Descent in the Accelerated Collatz Map}
\author{Wayne Brassem}

\begin{document}

% Custom commands
\NewDocumentCommand{\setN}{}{\mathbb{N}}				% Set of positive integers, excluding 0
\NewDocumentCommand{\setNo}{}{\mathbb{N}_0}				% Set of positive integers, including 0
\NewDocumentCommand{\setZ}{}{\mathbb{Z}}				% Set of all integers, excluding zero
\NewDocumentCommand{\setQ}{}{\mathbb{Q}}				% Set of all rational numbers
\newcommand{\chead}[1]{\multicolumn{1}{c}{#1}}

% Document commands which provide hyperlinks to OEIS sequences referenced herein
\NewDocumentCommand{\OEIS}{m}{\href{https://oeis.org/#1}{#1}}

\maketitle

\begin{abstract}

We study the Collatz map through the fully accelerated transformation
\[
f(x)=\frac{3x+1}{2^{v_2(3x+1)}},
\]
acting on the odd positive integers.

Finite accelerated trajectories are encoded by \emph{division words},
which record the sequence of 2-adic valuations encountered along an orbit.
Each admissible division word determines a canonical affine transformation
together with a unique congruence class modulo a power of two, yielding a
symbolic--affine representation of accelerated Collatz dynamics.

The admissible division words form a constrained symbolic language whose
admissibility is governed by the Beatty threshold sequence associated with
\(\log_2 3\). This Beatty grammar admits a canonical Sturmian return-word
decomposition determined by the continued-fraction expansion of \(\log_2 3\).
The resulting return blocks yield a spectral bound
\[
\lambda<3
\]
for the exponential growth rate of admissible symbolic configurations.
Comparing this spectral bound with the exponential dyadic refinement of the
associated congruence cylinders shows that symbolic proliferation is strictly
dominated by dyadic contraction

Consequently, the set of positive integers admitting an eventual descent
under the accelerated Collatz map has natural density one. The argument
does not exclude the existence of exceptional trajectories, but shows that
any such exceptions must belong to a residual set of natural density zero.

\end{abstract}

\newpage
\tableofcontents

\newpage
% ============================================================
\section{Introduction}
% ============================================================

The $3x+1$ problem, also known as the Collatz conjecture~\cite{Collatz1937}, concerns
the behavior of the map
\[
T(x) =
\begin{cases}
\displaystyle
    \phantom{3x} \frac{x}{2} & \text {if } x \text{ is even}, \\[8pt]
\displaystyle
    3x+1,                    & \text{if $x$ is odd}.
\end{cases}
\]
on the positive integers. It asserts that every positive integer eventually reaches $1$
under iteration of $T$. Despite its elementary formulation, the conjecture has resisted
proof for decades~\cite{Lagarias2010,Terras1976}. Extensive computational verification
confirms convergence for very large ranges~\cite{OliveiraESilva2010}, but does not
reveal the global structure governing trajectory behavior.

Recently, a major advance by Tao~\cite{Tao2019} established that almost all integers
eventually attain values smaller than their initial value under iteration of the Collatz
map in the sense of logarithmic density.

The present work develops a deterministic, arithmetic framework for analyzing typical
Collatz behavior based on a symbolic--affine encoding of trajectories. Rather than
studying individual trajectories directly, we encode the fully accelerated dynamics
symbolically using the map
\[
f(x)=\frac{3x+1}{2^{v_2(3x+1)}},
\]
which maps odd integers to odd integers by compressing each odd step into a single
transformation.

Each trajectory under $f$ is represented by a finite \emph{division word}, recording the
number of factors of $2$ removed (the 2-adic valuation) at each iterate. These words admit
an exact affine representation obtained by symbolic composition of the accelerated map:
each division word, \(\omega\), induces a map of the form
\[
F_\omega(x) = \frac{3^m x + c(\omega)}{2^{D(\omega)}},
\]
and determines a corresponding congruence class modulo $2^{D(\omega)}$. Thus division words
determine canonical arithmetic \emph{residue classes} in the integers, once realizability
is imposed.

Division words are \emph{admissible} if they arise from actual trajectories under the
accelerated map. Admissibility is governed by the Beatty~\cite{Graham1994} threshold sequence
\[
B(n) = \lfloor n \log_2 3 \rfloor + 1
\]
whose increment sequence is a Sturmian~\cite{Lothaire2002} word. Consequently the admissible
symbolic grammar inherits the combinatorial constraints of an irrational rotation rather than
arising from an ad hoc combinatorial construction.  This connection places the admissibility
problem within the well-developed theory of Beatty sequences and Sturmian symbolic dynamics.

Congruence classes allow the dynamics to be studied at the level of residue classes
rather than individual orbits. Convergent division words—those for which the associated
affine map is contractive—identify families of integers that decrease after a fixed
number of steps. The global problem is thereby reduced to understanding how these
families are distributed and how their associated densities accumulate.

We then show that the symbolic growth generated by the admissible grammar is strictly subcritical:
the number of admissible division words grows like \(\lambda^m\) with \(\lambda<3\), while the 
corresponding congruence cylinders undergo exponential dyadic refinement. The central mechanism
of the paper is the competition between these two exponential processes.

Admissible division words organize naturally into a forest of nested cylinders. Contractive
words define first-passage cylinders that remove positive measure from the residual set,
while non-contractive extensions form a shrinking residual forest.

This imbalance forces the residual forest to become exponentially sparse. Symbolic branching
cannot compensate for the increasingly fine arithmetic subdivision imposed by successive
congruence constraints.

The central result of the paper is that the admissible Beatty grammar has spectral growth
strictly below the rate of dyadic refinement. This reduces the density question to a comparison
of two exponential processes and implies that the residual forest has natural density zero.
Consequently, the set of integers admitting eventual descent under the accelerated Collatz map
has natural density one.

While this does not exclude the existence of exceptional trajectories, any such
exceptions must lie within a set of zero density and cannot arise from any scalable
arithmetic structure generated by the symbolic dynamics. Any counterexample to the Collatz
Conjecture, if it exists, must therefore avoid every contractive cylinder generated by the
affine grammar and remain confined to the residual forest.

\paragraph{Structure of the Paper.}
The argument proceeds as follows:
\begin{itemize}
\item Section~\ref{sec:division-counts} introduces division words and their total division counts.
\item Section~\ref{sec:affine-representation} develops the affine representation of division words and their congruence structure.
\item Section~\ref{sec:stepwise-realizability-admissibility} establishes the stepwise realizability of admissible division words and the unique residue lifting property.
\item Section~\ref{sec:cylinder-geometry} defines the geometry of contractive cylinders.
\item Section~\ref{sec:symbolic-growth} discusses competition between symbolic entropy and dyadic growth.
\item Section~\ref{sec:measure-collapse} establishes the collapse of measure for disjoint cylinders.
\item Section~\ref{sec:beatty-barrier} develops the Beatty grammar governing admissible symbolic growth and establishes the Beatty Barrier Theorem.
\item Section~\ref{sec:spectral-growth-bound} proves the spectral growth bound and combines it with the measure estimates to establish collapse of the residual forest.
\end{itemize}


% ============================================================
\section{Symbolic Encoding of Collatz Trajectories}
\label{sec:division-counts}
% ============================================================

% ------------------------------------------------------------
\subsection{Division Words}
% ------------------------------------------------------------

\begin{definition}[Division Word]
A division word of length $m \ge 1$ is a sequence
\[
\omega = (d_0, \dots, d_{m-1}),
\]
where each $d_i \in \mathbb{N}$.
\end{definition}

\begin{definition}[Admissible Division Word]
A division word
\[
\omega = (d_0,\dots,d_{m-1})
\]
is admissible if
\[
d_i \ge 1,
\qquad 0 \le i < m.
\]
\end{definition}

\begin{definition}[Realization]
An odd integer $x \in \mathbb{N}$ realizes the division word
\[
\omega=(d_0,\dots,d_{m-1})
\]
if
\[
d_i = v_2(3f^i(x)+1),
\qquad 0 \le i < m,
\]
where
\[
f(x)=\frac{3x+1}{2^{v_2(3x+1)}}.
\]
\end{definition}

\begin{remark}
Every odd integer \(x\) induces a unique infinite valuation sequence
\[
(v_2(3x+1),\, v_2(3f(x)+1),\, v_2(3f^2(x)+1),\dots),
\]
For every \(m\ge1\), its first \(m\) entries form the unique division
word of length \(m\) realized by x.
\end{remark}

% ------------------------------------------------------------
\subsection{Word Length and Total Division Count}
% ------------------------------------------------------------

\begin{definition}[Length]
The length of a division word
\[
\omega=(d_0,\ldots,d_{m-1})
\]
is
\[
|\omega|:=m.
\]
\end{definition}

\begin{definition}
The total division count of a word $\omega$ is
\[
D(\omega) := \sum_{i=0}^{m-1} d_i.
\]
\end{definition}

\begin{lemma}
For every admissible division word $\omega$,
\[
|\omega| \le D(\omega).
\]
\end{lemma}

\begin{proof}
Since $d_i\ge1$ for every $0\le i<|\omega|$,
\[
D(\omega)
=\sum_{i=0}^{|\omega|-1} d_i
\ge
\sum_{i=0}^{|\omega|-1}1
=
|\omega|.
\]
\end{proof}

\begin{remark}
The two basic numerical invariants of a division word are its length
\(|\omega|\), which records the number of accelerated Collatz steps, and
its total division count \(D(\omega)\), which records the cumulative
dyadic exponent. The interaction between these two quantities governs
both the affine dynamics and the symbolic geometry developed below.
\end{remark}


% ============================================================
\section{Affine Representation of Division Words}
\label{sec:affine-representation}
% ============================================================

\begin{definition}[Affine Map]
Each division word \(\omega\), whether or not it is realizable, determines
a formal affine map
\[
F_\omega : \mathbb Q \to \mathbb Q, \qquad
F_\omega(x)
=
\frac{3^{|\omega|}x+c(\omega)}
     {2^{D(\omega)}},
\]
obtained by symbolic composition of the accelerated Collatz steps encoded
by \(\omega\). Thus the word \(\omega\) determines its length \(|\omega|\),
its total division count \(D(\omega)\), and its affine constant \(c(\omega)\),
and hence determines the formal map \(F_\omega\).

The rational domain is used because an arbitrary formal division word need
not prescribe powers of \(2\) that divide the corresponding numerators
exactly. Thus the symbolic affine map is defined independently of
realizability. When a natural number \(x\) realizes \(\omega\), the prescribed
divisions are exact and \(F_\omega(x)\) agrees with the corresponding
iterate of the fully accelerated Collatz map.
\end{definition}

\begin{remark}
The quantity \(D(\omega)\) records the total exponent of \(2\) prescribed by
the symbolic segment encoded by \(\omega\). When \(\omega\) is realized, this
is the total exponent of \(2\) removed along the corresponding trajectory
segment.

The quantity \(2^{D(\omega)}\) is the denominator appearing in the formal
affine representation
\[
F_\omega(x)
=
\frac{3^{|\omega|}x+c(\omega)}{2^{D(\omega)}},
\]
and consequently determines the modulus of the associated formal
integrality condition
\[
3^{|\omega|}x+c(\omega)\equiv0 \pmod{2^{D(\omega)}}.\]
Exact realizability may impose additional congruence conditions beyond
this formal integrality constraint.

A division word is called \emph{realizable} if it is realized by at least
one odd integer.
\end{remark}

% ------------------------------------------------------------
\subsection{Concatenation and Prefix Structure}
% ------------------------------------------------------------

\begin{proposition}[Concatenation and Nested Affine Structure]
Let $\omega$ and $\eta$ be division words. If the concatenation
$\omega\eta$ denotes the word obtained by appending $\eta$ to $\omega$,
then the associated affine maps satisfy
\[
F_{\omega\eta} = F_\eta \circ F_\omega .
\]
Consequently, concatenation of division words is compatible with
composition of their affine maps: the affine map associated with
a concatenated word is exactly the composition of the affine maps
associated with its constituent words.

Furthermore, if $C_\omega$ denotes the set of natural numbers
realizing the word $\omega$, then
\[
C_{\omega\eta}\subseteq C_\omega .
\]
Thus extensions of division words correspond to nested realization
sets.
\end{proposition}

\begin{proof}
Write
\[
F_\omega(x) = \frac{3^{|\omega|}x+c(\omega)}{2^{D(\omega)}},
\]
and similarly
\[
F_\eta(x) = \frac{3^{|\eta|}x+c(\eta)}{2^{D(\eta)}}.
\]
Then
\[
F_\eta(F_\omega(x)) = \frac{ 3^{|\eta|} \left( \frac{3^{|\omega|} x+c(\omega)} {2^{D(\omega)}} \right)+c(\eta) } {2^{D(\eta)}}.
\]
Rearranging gives
\[
F_\eta(F_\omega(x)) = \frac{ 3^{|\omega|+|\eta|}x+3^{|\eta|} c(\omega)+2^{D(\omega)}c(\eta) } {2^{D(\omega)+D(\eta)}}.
\]
The denominator and ternary exponents add under concatenation, while
the affine constant combines according to the usual affine composition
rule:
\[
c(\omega\eta) = 3^{|\eta|} c(\omega)+2^{D(\omega)}c(\eta).
\]
Therefore this map is precisely the affine representation associated
with the concatenated word $\omega\eta$.

For the inclusion of realization sets, any integer realizing $\omega\eta$
must first realize the prefix $\omega$ before realizing the appended word
$\eta$. Hence every realization of $\omega\eta$ is also a realization of
$\omega$, and therefore
\[
C_{\omega\eta}\subseteq C_\omega .
\]
\end{proof}

\begin{lemma}[Affine Representation of a Realized Word]
If \(x\) realizes a division word \(\omega\), then
\[
F_\omega(x)=f^{|\omega|}(x). 
\]
In particular, although \(F_\omega\) is defined formally over \(\mathbb Q\)
for every division word, its value on a realization is an integer and
coincides with the corresponding finite segment of the fully accelerated
Collatz trajectory.
\end{lemma}

\begin{definition}[Contractive and Expansive Affine Maps]
Let $\omega$ be a realizable division word of length $m$. Since
\[
F_\omega(x)=f^{|\omega|}(x)
\]
for every realization of $\omega$, the affine map $F_\omega$ is called
\begin{itemize}
\item contractive if \(F_\omega(x)<x\) for every realization of $\omega$;
\item expansive if \(F_\omega(x)\ge x\) for every realization of $\omega$.
\end{itemize}
\end{definition}

\begin{remark}[No Dichotomy Assumed]
The preceding definitions do not assert that every realizable division word
is either contractive or expansive. At this stage it remains possible that
a realizable word $\omega$ has realizations $x$ and $y$ for which
\[
F_\omega(x)<x
\qquad\text{and}\qquad
F_\omega(y)\ge y.
\]
Such a word would satisfy neither definition.

Thus any later classification of realizable division words into contractive
and expansive classes requires an additional argument showing that the sign
of $F_\omega(x)-x$ is uniform over all realizations of a fixed word.
\end{remark}

\begin{remark}[Contractivity of a Finite Trajectory]
Contractivity refers only to the finite trajectory encoded by \(\omega\).
If \(F_\omega\) is contractive, then for every realization \(x\),
\[
F_\omega(x)=f^{|\omega|}(x)<x,
\]
This refers only to the finite trajectory encoded by \(\omega\); no conclusion
is drawn about the behavior of later iterates or eventual convergence to \(1\).
\end{remark}

% ------------------------------------------------------------
\subsection{Explicit Form of the Affine Constant}
% ------------------------------------------------------------

\begin{remark}[Division Words as a Dynamical Encoding]
Recall from Section~\ref{sec:division-counts} that a \emph{division word}
$\omega=(d_0,\dots,d_{m-1})$ prescribes a finite sequence of division counts.
When \(\omega\) is realized by a trajectory \(x_0,x_1,\dots\), its entries
are the exact \(2\)-adic valuations
\[
d_i = v_2(3 x_i + 1)
\]
occurring at each iterate of the fully accelerated Collatz map.

Division words retain the dynamical information needed to reconstruct the
affine representation. Each division word $\omega=(d_0,\dots,d_{m-1})$
determines the sequence of multiplicative factors of the linear coefficient
\[
r_i = \frac{3}{2^{d_i}},
\]
so that the multiplicative growth or decay contribution of each iterate is
preserved in the linear coefficient (with the affine term handled separately
via $c(\omega)$). This structure allows symbolic trajectories to be analyzed
through their affine representations and forms the foundation for the
realizability and descent results developed later.
\end{remark}

\begin{lemma}[Explicit Form of the Affine Constant]
Let $\omega = (d_0,\dots,d_{m-1})$ be a division word. Then the affine map
\[
F_\omega(x) = \frac{3^m x + c(\omega)}{2^{D(\omega)}}
\]
has the explicit form
\[
c(\omega)
=
\sum_{j=0}^{m-1}
3^{m-1-j}\cdot 2^{\sum_{i=0}^{j-1} d_i},
\]
where empty sums are defined as $0$.
\end{lemma}

\begin{proof}
Composing the symbolic steps prescribed by the division word gives
\[
x_{i+1}=\frac{3x_i+1}{2^{d_i}}.
\]

Repeated substitution yields
\[
x_m
=
\frac{3^m x_0}{2^{D(\omega)}}
+
\sum_{j=0}^{m-1}
\frac{
3^{m-1-j}
}{
2^{\sum_{i=j}^{m-1} d_i}
}.
\]
Multiplying by \(2^{D(\omega)}\) gives
\[
2^{D(\omega)}x_m
=
3^m x_0
+
\sum_{j=0}^{m-1}
3^{m-1-j}
2^{\sum_{i=0}^{j-1} d_i},
\]
which identifies the affine constant \(c(\omega)\).
\end{proof}

\begin{remark}
The affine constant \(c(\omega)\) records the accumulated contribution
of every additive term introduced during symbolic iteration. Earlier
contributions are multiplied by more subsequent factors of \(3\),
whereas later contributions inherit larger powers of \(2\) from the
preceding divisions.
\end{remark}

% ------------------------------------------------------------
\subsection{Dyadic-Ternary Balance}
% ------------------------------------------------------------

The symbolic dynamics of a division word are governed in part by the balance
between the cumulative ternary factor \(3^{|\omega|}\) and the cumulative
dyadic factor \(2^{D(\omega)}\). For a division word \(\omega\), the associated
affine map
\[
F_\omega(x)
=
\frac{3^{|\omega|}x+c(\omega)}
     {2^{D(\omega)}}
\]

has linear coefficient
\[
\lambda_\omega
=
\frac{3^{|\omega|}}{2^{D(\omega)}}.
\]

Thus
\[
2^{D(\omega)}<3^{|\omega|}
\]
corresponds to a linear coefficient greater than \(1\), whereas
\[
2^{D(\omega)}>3^{|\omega|}
\]
corresponds to a linear coefficient less than \(1\). The latter
condition should be distinguished from contractivity of the affine
map itself. For every nonempty division word \(\omega\), the affine
constant satisfies
\[
c(\omega)>0.
\]
Consequently, a linear coefficient less than \(1\) determines the
asymptotic slope of the affine map but does not by itself imply that
\[
F_\omega(x)<x
\]
for every realization \(x\). The linear coefficient itself, however,
admits a strict dichotomy for every nonempty division word. Indeed,
\[
2^{D(\omega)}\neq 3^{|\omega|}.
\]
If \(D(\omega)>0\), then \(2^{D(\omega)}\) is even whereas
\(3^{|\omega|}\) is odd. If \(D(\omega)=0\), then
\(2^{D(\omega)}=1\), while \(3^{|\omega|}>1\) because
\(|\omega|>0\). Hence exactly one of
\[
2^{D(\omega)}<3^{|\omega|}
\qquad\text{or}\qquad
2^{D(\omega)}>3^{|\omega|}
\]
holds, and therefore
\[
\lambda_\omega>1
\qquad\text{or}\qquad
\lambda_\omega<1,\]
respectively.

This slope dichotomy should not be confused with a dichotomy between
expansive and contractive affine maps. The positive affine constant can
affect the sign of \(F_\omega(x)-x\), so a subunit linear coefficient does
not alone establish contractivity on the realizations of \(\omega\).

The boundary between the two slope regimes is described by the dyadic
threshold associated with a word length \(m\),
\[
B(m)=\lfloor m\log_2 3\rfloor+1.
\]

Equivalently, \(B(m)\) is the least natural number \(k\) satisfying
\[
2^k>3^m.
\]

The sequence \(B(m)\) is OEIS sequence \OEIS{A020914}. Thus,
for a division word \(\omega\),
\[
B(|\omega|)\le D(\omega)
\quad\Longleftrightarrow\quad
2^{D(\omega)}>3^{|\omega|}.\]

Therefore \(B(|\omega|)\) is exactly the minimum total division count for
which the linear coefficient of \(F_\omega\) is strictly less than \(1\).

% ------------------------------------------------------------
\subsection{Example}
% ------------------------------------------------------------

Consider division words of length \(m=4\). The dyadic threshold is
\[
B(4)=A020914(4)=7,
\]
so a total division count \(D(\omega)=7\) is the minimum for which the
linear coefficient of the associated affine map is strictly less than \(1\). 
Consider the formal division word
\[
\omega=(1,1,1,4),
\]
for which
\[
D(\omega)=1+1+1+4=7.
\]
Its affine constant is
\[
c(\omega)=65,
\]
and hence
\[
F_\omega(x) = \frac{3^4x+c(\omega)}{2^7} = \frac{81x+65}{128}.
\]
The congruence
\[
81x+65\equiv0\pmod{128}
\]
ensures that the formal affine expression is integral. It determines the
residue class
\[
x\equiv15\pmod{128}.
\]
This condition alone, however, does not guarantee that \(\omega\) is
realized by the fully accelerated Collatz map, because realization requires
each prescribed division count to be the exact \(2\)-adic valuation at the
corresponding step.

For example, \(x=15\) satisfies the congruence modulo \(128\), but its
accelerated orbit begins
\[
15\longmapsto23\longmapsto35\longmapsto53\longmapsto5.
\]
At the fourth step,
\[
3\cdot53+1=160=2^5\cdot5,
\]
so the corresponding valuation is \(5\), not \(4\). The formal division
prescribed by \(\omega\) instead gives
\[
\frac{3\cdot53+1}{2^4}=10,
\]
which is integral but even. Thus \(15\) satisfies the formal integrality
condition without realizing the division word \((1,1,1,4)\). Exact
realization in particular requires the quotient after the prescribed four
divisions to be odd. This additional parity condition is
\[
\frac{81x+65}{128}\equiv1\pmod2,
\]
or equivalently,
\[
81x+65\equiv128\pmod{256}.
\]
Solving this strengthened congruence gives
\[
x\equiv143\pmod{256}.
\]
For \(x=143\), the accelerated orbit is
\[
143\longmapsto215\longmapsto323\longmapsto485\longmapsto91,
\]
and the successive \(2\)-adic valuations are exactly
\[
(1,1,1,4).
\]
Thus \(143\) genuinely realizes the division word \(\omega\).

\medskip

\noindent
\textbf{Interpretation}

The example distinguishes two related arithmetic conditions. The
congruence modulo \(2^{D(\omega)}\) guarantees that the symbolic affine
composition is integral, while exact realization requires the prescribed
powers of \(2\) to be maximal at every step.

In this example, imposing the necessary final-parity condition contributes
one additional congruence constraint, refining
\[
x\equiv15\pmod{128}
\]
to
\[
x\equiv143\pmod{256}.
\]
Thus the affine representation converts symbolic trajectory data into
explicit congruence conditions, but integrality of the formal affine map
must be distinguished from exact realizability under the fully accelerated
Collatz map. The following sections develop the additional conditions
required to characterize realizable division words and their associated
cylinders.


% ============================================================
\section{Stepwise Realizability and Admissibility}
\label{sec:stepwise-realizability-admissibility}
% ============================================================

% ------------------------------------------------------------
\subsection{Stepwise Realizability}
% ------------------------------------------------------------

\begin{theorem}[Affine Stepwise Identity]
\label{thm:affine-stepwise-identity}
Let $\omega = (d_0,\dots,d_{m-1})$ be a division word, and define $c(\omega)$ by
\[
c(\omega)
=
\sum_{j=0}^{m-1}
3^{m-1-j}\cdot 2^{\sum_{i=0}^{j-1} d_i}.
\]

% Suppose the successive divisions by \(2^{d_i}\) are all exact, so that
Suppose
\[
v_2(3x+1) = d_i, \qquad 0 \le i < m.
\]
% f(x)=\frac{3x+1}{2^{v_2(3x+1)}},
% v_2(3f^j(x)+1)=d_j, \qquad 0\le j<m.
Equivalently,
\[
x_{i+1} = \frac{3x_i+1}{2^{d_i}}, \qquad 0\le i<m.
\]
Then
\[
2^{D(\omega)} x_m = 3^m x_0 + c(\omega).
\]
\end{theorem}

\begin{corollary}[Stepwise Integrality Criterion]
The sequence $(x_0,\dots,x_m)$ is an integer trajectory
under the accelerated Collatz relations if and only if
\[
3^m x_0 + c(\omega) \equiv 0 \pmod{2^{D(\omega)}}.
\]
\end{corollary}

\begin{remark}[Structural Interpretation]
Each term in $c(\omega)$ records the propagated contribution of a
single additive ``+1'' created during iteration. Earlier contributions
accumulate additional powers of $3$ through forward propagation and
additional powers of $2$ through delayed subsequent divisions by powers
of two, producing a fully time-ordered encoding of the trajectory in a
single affine constant.
\end{remark}

% ------------------------------------------------------------
\subsection{Division Word Valuation}
% ------------------------------------------------------------

\begin{lemma}[Exact Valuation Congruence]
\label{lem:exact-valuation-congruence}
Let \(x\) be an odd integer and let \(d\ge1\).

Then
\[
v_2(3x+1)=d
\]
if and only if
\[
3x+1\equiv 2^d \pmod{2^{d+1}}.
\]
Equivalently,
\[
3x+1 = 2^d(2k+1)
\]
for some integer \(k\).
\end{lemma}

\begin{proof}
Suppose first that
\[
v_2(3x+1)=d.
\]
Then by definition,
\[
3x+1 = 2^d y,
\]
where \(y\) is odd. Writing
\[
y=2k+1,
\]
gives
\[
3x+1 = 2^d(2k+1) = 2^{d+1}k+2^d,
\]
and therefore
\[
3x+1 \equiv 2^d \pmod{2^{d+1}}.
\]

Conversely, suppose that
\[
3x+1 \equiv 2^d \pmod{2^{d+1}}.
\]
Then there exists an integer \(k\) such that
\[
3x+1 = 2^{d+1}k+2^d = 2^d(2k+1).
\]
Since \(2k+1\) is odd, exactly \(d\) factors of \(2\) divide \(3x+1\). Hence
\[
v_2(3x+1)=d.
\]
\end{proof}

\begin{corollary}[Local Valuation Refinement]
\label{cor:local-refinement}
Let
\[
\omega=(d_0,\dots,d_{m-1})
\]
be an admissible division word and let
\[
\omega_j=(d_0,\dots,d_{j-1})
\]
denote its prefixes. For each \(j\), the condition
\[
v_2(3f^j(x)+1)=d_j
\]
is equivalent to the congruence
\[
3f^j(x)+1 \equiv 2^{d_j} \pmod{2^{d_j+1}}.
\]

Consequently, each additional valuation constraint refines the realization
set of the prefix \(\omega_j\) to a smaller realization set associated with
the extended prefix
\[
\omega_{j+1}.
\]

Thus admissible division words generate a nested sequence of arithmetic
constraints, and extensions of division words correspond to successive
refinements of residue classes.
\end{corollary}

\begin{remark}
The corollary formalizes the local nature of admissibility. Exact
valuation constraints are imposed one step at a time and depend only on
the current accelerated iterate. Consequently, realizability propagates
forward recursively through the valuation sequence rather than being
determined solely by a terminal congruence condition.
\end{remark}

\begin{theorem}[Residue Lifting Theorem]
\label{thm:residue-lifting-theorem}
Let
\[
\omega=(d_0,\dots,d_{m-1})
\]
be an admissible division word.  Then there exists a unique residue class
\[
r_\omega \pmod{2^{D(\omega)}}
\]
consisting precisely of the integers realizing \(\omega\).  Moreover, if
\[
\omega_j=(d_0,\dots,d_{j-1})
\]
denotes the prefix of length \(j\), then the corresponding residue classes satisfy
\[
r_{\omega_{j+1}}
\equiv
r_{\omega_j}
\pmod{2^{D(\omega_j)}}.
\]

Thus each additional valuation constraint uniquely refines the residue
class of its prefix.
\end{theorem}

\begin{proof}
We proceed by induction on the length \(m\) of the division word.
For \(m=1\), the word consists of a single valuation
\[
\omega=(d_0).
\]
By Lemma~\ref{lem:exact-valuation-congruence},
\[
v_2(3x+1)=d_0
\]
if and only if
\[
3x+1
\equiv
2^{d_0}
\pmod{2^{d_0+1}}.
\]

Since \(3\) is invertible modulo every power of two, this congruence has
a unique solution modulo \(2^{d_0+1}\). Thus the realization set of
\(\omega\) is a single residue class.

Assume now that the result holds for all admissible words of length \(m\),
and let
\[
\omega'=(d_0,\dots,d_{m})
\]
be an admissible word of length \(m+1\).  Let
\[
\omega=(d_0,\dots,d_{m-1})
\]
be its length-\(m\) prefix. By the induction hypothesis, the integers
realizing \(\omega\) form a unique residue class
\[
r_\omega
\pmod{2^{D(\omega)}}.
\]

The additional valuation condition
\[
v_2(3f^{m}(x)+1)=d_m
\]
is equivalent, by
Lemma~\ref{lem:exact-valuation-congruence}, to
\[
3f^{m}(x)+1
\equiv
2^{d_m}
\pmod{2^{d_m+1}}.
\]

Using the affine representation of the prefix,
\[
f^{m}(x) = \frac{3^{m}x+c(\omega)}{2^{D(\omega)}},
\]
this becomes a linear congruence condition on \(x\) modulo
\[
2^{D(\omega)+d_m} = 2^{D(\omega')}.
\]

Since \(3^m\) is odd, it is invertible modulo every power of two.  This congruence
determines a unique refinement of the residue class
\[
r_\omega
\pmod{2^{D(\omega)}}
\]
to a residue class
\[
r_{\omega'}
\pmod{2^{D(\omega')}}.
\]

Therefore every admissible extension uniquely refines the residue class
of its prefix.  The result follows by induction.
\end{proof}

% ------------------------------------------------------------
\subsection{Affine Realization}
% ------------------------------------------------------------

\begin{theorem}[Affine Realization Theorem]
\label{thm:affine-realization}
Let $\omega$ be a division word.

An integer $x$ realizes $\omega$ if and only if
\[
x\equiv r_\omega\pmod{2^{D(\omega)}}.
\]

\end{theorem}

\begin{proof}
Assume that
\[
x\equiv r_\omega \pmod{2^{D(\omega)}}.
\]

By the Exact Valuation Congruence Lemma \ref{lem:exact-valuation-congruence},
the residue class associated with the full word $\omega$ is contained in the residue
class of every prefix
\[
\omega_j=(d_0,\dots,d_{j-1}).
\]
Therefore, for every prefix length $j$,
\[
x\equiv r_{\omega_j}\pmod{2^{D(\omega_j)}}.
\]

We prove inductively that the successive iterates generated from $x$ realize
the prescribed valuations.  For $j=0$, the prefix condition gives
\[
x\equiv r_{(d_0)}\pmod{2^{d_0}},
\]
and therefore
\[
2^{d_0}\mid 3x+1.
\]
Because the residue class is defined using the exact division count $d_0$,
the next higher power does not divide:
\[
2^{d_0+1}\nmid 3x+1.
\]
Hence
\[
v_2(3x+1)=d_0.
\]
Therefore the first accelerated iterate is
\[
x_1=\frac{3x_0+1}{2^{d_0}},
\]
and is an integer.  Now assume that the first $j$ steps have been realized, so that
\[
x_j=f^j(x)
\]
is the integer determined by the prefix
\[
(d_0,\dots,d_{j-1}).
\]
The prefix congruence condition for
\[
\omega_{j+1}=(d_0,\dots,d_j)
\]
implies that
\[
2^{d_j}\mid 3x_j+1.
\]
Again, exactness of the prefix residue condition excludes the higher power:
\[
2^{d_j+1}\nmid 3x_j+1.
\]
Consequently,
\[
v_2(3x_j+1)=d_j,
\]
and the next iterate
\[
x_{j+1}=\frac{3x_j+1}{2^{d_j}}
\]
is an integer.  By induction, every step satisfies the prescribed valuation
constraint, and therefore
\[
v_2(3f^j(x)+1)=d_j, \qquad 0\le j<m.
\]
Hence $x$ realizes the division word $\omega$.
\end{proof}

% ------------------------------------------------------------
\subsection{Admissibility of Division Words}
% ------------------------------------------------------------

\begin{definition}[Admissible Division Word]
\label{def:admissible-division-word}
A division word is \emph{admissible} if it is realizable.
\end{definition}

\begin{remark}
Equivalently, a division word \(\omega\) is admissible if its induced affine
residue class contains an integer whose successive accelerated iterates satisfy
\[
v_2(3x_i+1)=d_i,
\qquad 0\le i<m.
\]
Thus admissibility is precisely the condition that a formal symbolic word
occurs as an actual finite valuation sequence of the accelerated Collatz map.
\end{remark}

% ------------------------------------------------------------
\subsection{Residue Class Representation of Division Words}
% ------------------------------------------------------------

\begin{lemma}[Prefix Nesting of Realization Sets]
\label{lem:prefix-uniqueness}
Let $\omega,\omega'$ be admissible division words.

If the realization sets of $\omega$ and $\omega'$ intersect, then one word
is a prefix of the other.
\end{lemma}

\begin{proof}
Every odd integer \(x\) determines a unique infinite valuation sequence
\[
(v_2(3x+1),\, v_2(3f(x)+1),\, v_2(3f^2(x)+1),\dots).
\]

If \(x\) lies in the realization sets of both \(\omega\) and \(\omega'\),
then both words must agree with the initial segment of this valuation sequence.

Therefore the two words cannot diverge at any position before the shorter
word terminates. Hence one word is a prefix of the other.
\end{proof}


% ============================================================
\section{Cylinder Geometry of the Affine Grammar}
\label{sec:cylinder-geometry}
% ============================================================

The admissible division-word grammar is organized by
2-adic residue classes induced by admissible division words.

These residue classes provide an arithmetic realization of the
symbolic dynamics and naturally suggest a geometric interpretation
in terms of nested congruence classes, developed below.

% ------------------------------------------------------------
\subsection{Cylinder of a Division Word}
% ------------------------------------------------------------

\begin{definition}[Cylinder]
Let $\omega$ be an admissible division word with associated residue
class
\[
r_\omega \pmod{2^{D(\omega)}}.
\]

The cylinder associated with $\omega$ is
\[
C_\omega := \{x \in \mathbb N : x \equiv r_\omega \pmod{2^{D(\omega)}}\}.
\]
\end{definition}

% ------------------------------------------------------------
\subsection{Prefix Structure and Forest Geometry of Cylinders}
% ------------------------------------------------------------

\begin{lemma}[Prefix Nesting of Cylinders]
\label{lem:cylinder-prefix-nesting}
Let $\omega$ and $\omega'$ be admissible division words.

If
\[
C_\omega\cap C_{\omega'}\neq\varnothing,
\]
then one of the words is a prefix of the other.
\end{lemma}

\begin{proof}
Suppose
\[
x\in C_\omega\cap C_{\omega'}.
\]

Then $x$ realizes both division words. By the Affine Realization
Theorem~\ref{thm:affine-realization}, $C_\omega$ and $C_{\omega'}$
are precisely the realization sets of $\omega$ and $\omega'$.
The result therefore follows immediately from
Lemma~\ref{lem:prefix-uniqueness}.
\end{proof}

\begin{theorem}[Cylinder Inclusion Criterion]
\label{thm:cylinder-inclusion}
Let $\omega$ and $\eta$ be admissible division words.

Then
\[
C_\eta\subseteq C_\omega
\]
if and only if
\[
\omega
\]
is a prefix of
\[
\eta.
\]
\end{theorem}

\begin{proof}
Suppose first that $\omega$ is a prefix of $\eta$.
Any integer realizing $\eta$ must realize every prefix of $\eta$,
including $\omega$. Therefore
\[
C_\eta\subseteq C_\omega.
\]

Conversely, suppose
\[
C_\eta\subseteq C_\omega.
\]

Let
\[
x\in C_\eta.
\]

Then
\[
x\in C_\eta\cap C_\omega.
\]

By Lemma~\ref{lem:cylinder-prefix-nesting}, one word is a prefix
of the other. If $\eta$ were a prefix of $\omega$, then applying
the forward direction already established would imply
\[
C_\omega\subseteq C_\eta.
\]

Together with
\[
C_\eta\subseteq C_\omega,
\]
this would force
\[
C_\eta=C_\omega,
\]
Since each cylinder is defined by a unique residue class modulo
\(2^{D(\omega)}\), equality of cylinders implies equality of moduli.  Thus
\[
D(\eta)=D(\omega).
\]
Because one word is a prefix of the other and the two words have equal
depth, they must coincide:
\[
\eta=\omega.
\]
Therefore $\omega$ must be a prefix of $\eta$.
\end{proof}

\begin{corollary}[Disjointness of Sibling Cylinders]
\label{cor:sibling-disjointness}
If neither admissible word is a prefix of the other, then
\[
C_\omega\cap C_\eta=\varnothing.
\]
\end{corollary}

\begin{proof}
If the cylinders intersected, then by
Lemma~\ref{lem:cylinder-prefix-nesting},
one word would be a prefix of the other, contradicting the assumption.
\end{proof}

\begin{corollary}[Disjointness at Equal Depth]
\label{cor:equal-depth-disjoint}
Distinct admissible division words of equal length define disjoint
cylinders.
\end{corollary}

\begin{proof}
If two equal-length words intersected, then by
Lemma~\ref{lem:cylinder-prefix-nesting},
one would be a prefix of the other.
Equal lengths force equality of the words, contradicting distinctness.
\end{proof}

\begin{corollary}[Symbolic-Arithmetic Correspondence]
\[
\omega\preceq\eta
\iff
C_\eta\subseteq C_\omega.
\]
\end{corollary}

\begin{proof}
This is exactly the content of
Theorem~\ref{thm:cylinder-inclusion}.
\end{proof}

\begin{remark}
The previous results identify symbolic ancestry with arithmetic
containment.  Thus admissible division words organize into a collection
of nested 2-adic cylinders whose inclusion relations coincide exactly
with the prefix order on division words.

The admissible symbolic grammar therefore forms a disjoint \emph{forest}
of prefix trees rather than a single global tree. Each node corresponds
to an admissible one-step refinement of a cylinder, obtained by
extending the associated division word by one admissible symbol.  Each
descendant cylinder is obtained by refining the associated congruence
class modulo a higher power of two.
\end{remark}

% ------------------------------------------------------------
\subsection{Contractive Cylinders}
% ------------------------------------------------------------

\begin{definition}[Contractive Cylinder]
A cylinder \(C_\omega\) is called linearly contractive if its
associated affine map
\[
F_\omega(x) = \frac{3^{|\omega|}x+c(\omega)}{2^{D(\omega)}}
\]
has contractive linear coefficient:
\[
\frac{3^{|\omega|}}{2^{D(\omega)}}<1.
\]
Equivalently,
\[
D(\omega)>|\omega|\log_2 3.
\]
\end{definition}

% ------------------------------------------------------------
\subsection{Cylinder Contraction Threshold}
% ------------------------------------------------------------

\begin{lemma}[Minimal Contraction Threshold]
\label{lem:minimal-contraction-threshold}
Let $\omega$ be an admissible division word of length $m$. If the
associated affine map
\[
F_\omega(x) = \frac{3^m x+c(\omega)}{2^{D(\omega)}}
\]
is contractive in its linear coefficient, then
\[
D(\omega)\ge A020914(m),
\]
where
\[
B(m) = A020914(m)=\lfloor m\log_2 3\rfloor + 1
\]
is the minimal integer \(k\) satisfying
\[
2^k>3^m.
\]
\end{lemma}

\begin{proof}
Contractivity requires
\[
\frac{3^m}{2^{D(\omega)}}<1.
\]
which is equivalent to
\[
2^{D(\omega)}>3^m.
\]

The minimal integer satisfying this inequality is exactly
\(A020914(m)\), so
\[
D(\omega) \ge B(m) = A020914(m).
\]
\end{proof}

\begin{remark}
The Beatty threshold sequence \OEIS{A020914} governs the growth/decay balance along
admissible branches within every cylinder of the forest.
\end{remark}


% ============================================================
\section{Symbolic Growth of the Cylinder Forest}
\label{sec:symbolic-growth}
% ============================================================

The cylinder forest introduced in the previous section exhibits two
competing exponential effects:

\begin{enumerate}
\item combinatorial growth in the number of admissible cylinders; and
\item dyadic contraction in the size of individual cylinders.
\end{enumerate}

The first of these is purely symbolic and is quantified below.

% ------------------------------------------------------------
\subsection{Symbolic Complexity}
% ------------------------------------------------------------

\begin{definition}[Symbolic complexity]
Let \(N(m)\) denote the number of admissible division words of length \(m\).
Equivalently, \(N(m)\) is the number of admissible cylinders appearing
at depth \(m\) in the symbolic forest.
\end{definition}

\begin{remark}
The quantity \(N(m)\) measures the combinatorial complexity of the admissible
grammar at depth \(m\). It counts the number of distinct admissible division
words, or equivalently the number of admissible cylinders obtained after \(m\)
levels of symbolic refinement.
\end{remark}

% ------------------------------------------------------------
\subsection{Symbolic Growth Rate}
% ------------------------------------------------------------

\begin{definition}[Symbolic Growth Rate]
Define the symbolic growth rate of the admissible grammar by
\[
\lambda := \limsup_{m\to\infty} N(m)^{1/m}.
\]
\end{definition}

\begin{remark}
The quantity \(\lambda\) is the exponential branching rate of the
admissible cylinder forest. It measures the asymptotic combinatorial
complexity of admissible symbolic extensions.
\end{remark}

\begin{lemma}[Uniform exponential envelope]
\label{lem:uniform-envelope}
For every \(\varepsilon>0\), there exists \(m_0\) such that
\[
N(m) \le (\lambda+\varepsilon)^m
\]
for all \(m\ge m_0\).
\end{lemma}

\begin{proof}
This is an immediate consequence of the definition of
\[
\lambda := \limsup_{m\to\infty} N(m)^{1/m}.
\]
\end{proof}

% ------------------------------------------------------------
\subsection{Entropy Versus Dyadic Contraction}
% ------------------------------------------------------------

Each admissible cylinder is determined by a congruence class modulo
\[
2^{D(\omega)}
\]
Consequently, deeper cylinders correspond to congruence classes modulo
progressively larger powers of two. Since
\[
D(\omega) \ge B(m)
\]
every cylinder at depth \(m\) is defined modulo at least
\[
2^{B(m)}
\]
where
\[
B(m) = A020914(m) = \lfloor m\log_2 3\rfloor + 1
\]
therefore
\[
2^{-D(\omega)} \le 2^{-B(m)} < 3^{-m}.
\]

Two independent quantities govern the structure of the depth-\(m\) cylinder
forest: the number \(N(m)\) of admissible cylinders and the modulus
\(2^{D(\omega)}\) associated with each cylinder.

The following section assigns a natural density to each cylinder.
Since cylinders at equal depth are pairwise disjoint
(Corollary~\ref{cor:equal-depth-disjoint}), the residual density of
the depth-\(m\) cylinder forest is bounded above by the product of the
number of admissible cylinders and the maximal density of any one cylinder.


% ============================================================
\section{Measure Collapse of the Residual Forest}
\label{sec:measure-collapse}
% ============================================================

The admissible forest constructed in the previous sections induces a
nested family of subsets of the odd integers. These sets encode the
integers whose symbolic trajectories remain admissible to increasing
depths.

The central question is whether these residual sets \(S_m\) retain
positive natural density under repeated symbolic refinement.

% ------------------------------------------------------------
\subsection{Residual Cylinder Sets}
% ------------------------------------------------------------

\begin{definition}[Depth-$m$ Residual Set]
Let
\[
S_m := \bigcup_{|\omega|=m} C_\omega
\]
where the union ranges over all admissible division words of length \(m\).
\end{definition}

\begin{remark}
The set \(S_m\) consists of integers whose valuation trajectories admit
an admissible division-word prefix of length \(m\). Equivalently, these
are precisely the trajectories that remain unresolved after \(m\)
refinement steps.
\end{remark}

% ------------------------------------------------------------
\subsection{First-Passage Layers}
% ------------------------------------------------------------

\begin{definition}[First-Passage Layer]
Define the first-passage layer at depth \(m\) by
\[
F_m := S_m \setminus S_{m+1}.
\]
\end{definition}

\begin{remark}
The set \(F_m\) consists of integers whose symbolic trajectories
admit an admissible prefix of length \(m\), but admit no admissible
extension to depth \(m+1\).
\end{remark}

\begin{lemma}[Disjoint residual decomposition]
For every \(m\ge1\),
\[
S_m = F_m \sqcup S_{m+1},
\]
where \(\sqcup\) denotes disjoint union.
\end{lemma}

\begin{proof}
By definition,
\[
F_m = S_m \setminus S_{m+1}.
\]
Hence every element of \(S_m\) either:
\begin{itemize}
\item remains unresolved to depth \(m+1\), or
\item exits the residual forest through the first-passage layer \(F_m\)
\end{itemize}

These two possibilities are mutually exclusive, giving the disjoint decomposition.
\end{proof}

\begin{remark}[Residual decomposition]
The decomposition
\[
S_m = F_m \sqcup S_{m+1}
\]
shows that symbolic refinement partitions the residual set into two
disjoint components: trajectories that exit the residual forest at
depth \(m\), and trajectories that remain unresolved to the next level.
No trajectories are created or duplicated under refinement.
\end{remark}

\begin{corollary}[Global first-passage decomposition]
\label{cor:global-decomposition}
The initial residual set admits the disjoint decomposition
\[
S_1 = \left(\bigsqcup_{m\ge1} F_m\right) \sqcup S_\infty,
\]
where
\[
S_\infty = \bigcap_{m\ge1} S_m
\]
is the infinite residual set.
\end{corollary}

\begin{proof}
Iterating the decomposition
\[
S_m = F_m \sqcup S_{m+1}
\]
gives
\[
S_1 = F_1 \sqcup F_2 \sqcup \cdots \sqcup F_m \sqcup S_{m+1}.
\]

Passing to the limit \(m\to\infty\) yields
\[
S_1 = \left(\bigsqcup_{m\ge1}F_m\right) \sqcup S_\infty.
\qedhere
\]
\end{proof}

\begin{remark}[Symbolic conservation law]
The global decomposition shows that every residual trajectory
either:
\begin{itemize}
\item eventually exits into a first-passage layer, or
\item survives indefinitely in \(S_\infty\).
\end{itemize}

No symbolic mass is lost or duplicated under refinement; it is
partitioned disjointly between finite first-passage escape layers
and the infinite residual set.
\end{remark}

% ------------------------------------------------------------
\subsection{Cylinder Measure}
% ------------------------------------------------------------

\begin{lemma}[Cylinder Density]
\label{lem:cylinder-density}
Let \(C_\omega\) be an admissible cylinder associated with the residue
class
\[
r_\omega \pmod{2^{D(\omega)}}.
\]
Then
\[
\mu(C_\omega)=2^{-D(\omega)},
\]
where \(\mu\) denotes natural density on the odd positive integers.
\end{lemma}

\begin{proof}
By the Affine Realization Theorem~\ref{thm:affine-realization},
\(C_\omega\) is exactly one residue class modulo \(2^{D(\omega)}\).
A residue classes modulo \(2^{D(\omega)}\) has natural density
\(2^{-D(\omega)}\). Therefore
\[
\mu(C_\omega)=2^{-D(\omega)}.
\]
\end{proof}

\begin{remark}
The density of a cylinder depends only on its modulus \(2^{D(\omega)}\),
not on the particular residue \(r_\omega\). Thus deeper cylinders are
exponentially smaller, with each additional power of two halving the
natural density.
\end{remark}

\begin{lemma}[Disjointness at fixed depth]
Distinct admissible cylinders of equal depth are disjoint.
\end{lemma}

\begin{proof}
This follows from the prefix-nesting property established in
Section~\ref{sec:cylinder-geometry}.
\end{proof}

\begin{corollary}[Measure decomposition]
\label{cor:measure-decomposition}
For every \(m\),
\[
\mu(S_m) = \sum_{|\omega|=m}2^{-D(\omega)}.
\]
\end{corollary}

\begin{remark}[Conservation under symbolic refinement]
The first-passage decomposition partitions the initial residual set
into disjoint escape layers together with the infinite residual set.
Thus dyadic mass is neither created nor lost under refinement:
it is redistributed between progressively finer symbolic layers.
\end{remark}

% ------------------------------------------------------------
\subsection{Nested residual structure}
% ------------------------------------------------------------

\begin{definition}[Infinite Residual Set]
Define the infinite residual set by
\[
S_\infty := \bigcap_{m=1}^{\infty} S_m.
\]
\end{definition}

\begin{remark}
The set \(S_\infty\) consists of integers whose symbolic trajectories
admit admissible extensions at every finite depth, equivalently,
integers which never exit the residual forest under successive refinement.
\end{remark}

\begin{theorem}[Continuity of measure under nested refinement]
\label{thm:nested-measure}
Since
\[
S_{m+1}\subseteq S_m,
\]
the measures satisfy
\[
\mu(S_\infty) = \lim_{m\to\infty}\mu(S_m).
\]
\end{theorem}

\begin{proof}
This follows from continuity from above for natural density on
nested measurable residue-class unions.
\end{proof}

% ------------------------------------------------------------
\subsection{Residual Density Bound}
% ------------------------------------------------------------

\begin{theorem}[Residual Density Bound]
\label{thm:residual-density-bound}

For every depth \(m\),
\[
\mu(S_m) = \sum_{|\omega|=m}2^{-D(\omega)} \le N(m)\,2^{-B(m)},
\]
where
\[
B(m)=A020914(m) = \lfloor m\log_2 3\rfloor+1.
\]
\end{theorem}

\begin{proof}
By Corollary~\ref{cor:measure-decomposition},
\[
\mu(S_m) = \sum_{|\omega|=m} 2^{-D(\omega)}.
\]
By Lemma~\ref{lem:minimal-contraction-threshold},
\[
D(\omega)\ge B(m)
\]
for every admissible division word of length \(m\). Therefore
\[
2^{-D(\omega)} \le 2^{-B(m)}.
\]
Since the sum ranges over \(N(m)\) admissible division words of length \(m\),
\[
\mu(S_m) \le N(m)\,2^{-B(m)},
\]
as claimed.
\end{proof}

\begin{theorem}[Residual Density Collapse]
\label{thm:residual-density-collapse}
Suppose the symbolic growth rate satisfies
\[
\lambda<3,
\]
where
\[
\lambda = \limsup_{m\to\infty} N(m)^{1/m}.
\]
Then
\[
\mu(S_m) \longrightarrow 0.
\]
Consequently,
\[
\mu(S_\infty) = 0.
\]
\end{theorem}

\begin{proof}
Let
\[
\varepsilon > 0
\]
be chosen so that
\[
\lambda + \varepsilon < 3.
\]
By Lemma~\ref{lem:uniform-envelope}, there exists
\(m_0\) such that
\[
N(m) \le (\lambda+\varepsilon)^m
\] 
for every \(m\ge m_0\).  Combining this estimate with
Theorem~\ref{thm:residual-density-bound} gives
\[
\mu(S_m) \le (\lambda+\varepsilon)^m \,2^{-B(m)}.
\]
Since
\[
2^{B(m)} > 3^m,
\]
we have
\[
2^{-B(m)} < 3^{-m}.
\]
Hence
\[ \mu(S_m) < \left( \frac{\lambda+\varepsilon}{3} \right)^m.
\]
Because
\[
\frac{\lambda+\varepsilon} {3} < 1,
\]
the right-hand side converges exponentially to zero. Therefore
\[
\mu(S_m) \to 0.
\]
Finally, Theorem~\ref{thm:nested-measure} implies
\[
\mu(S_\infty) = \lim_{m\to\infty}\mu(S_m) = 0.
\]
\end{proof}

% ------------------------------------------------------------
\subsection{First-Passage Completeness}
% ------------------------------------------------------------

\begin{theorem}[First-Passage Completeness]
\label{thm:first-passage-completeness}
Suppose
\[
\lambda < 3.
\]
Then
\[
\mu(S_1) = \sum_{m=1}^{\infty} \mu(F_m).
\]
In particular, if the initial residual set is normalized so that
\[
\mu(S_1) = 1,
\]
then
\[
\sum_{m=1}^{\infty} \mu(F_m) = 1.
\]
Thus the set of integers which remain in the residual forest
indefinitely has natural density zero.
\end{theorem}

\begin{proof}
By Corollary~\ref{cor:global-decomposition},
\[
S_1 = \left( \bigsqcup_{m\ge1}F_m \right) \sqcup S_\infty.
\]
Since the union is disjoint,
\[
\mu(S_1) = \sum_{m\ge1}\mu(F_m) + \mu(S_\infty).
\]
By Theorem~\ref{thm:residual-density-collapse},
\[
\mu(S_\infty) = 0.
\]
Therefore
\[
\mu(S_1) = \sum_{m\ge1} \mu(F_m).
\]
If the measure is normalized so that \(\mu(S_1)=1\),
the stated identity follows immediately.
\end{proof}

% ------------------------------------------------------------
\subsection{Interpretation}
% ------------------------------------------------------------

The preceding results reduce the density-one problem for accelerated Collatz
trajectories to a purely symbolic question. Every measure-theoretic aspect of
the argument is now contained in Theorem~\ref{thm:residual-density-collapse}.
Consequently, it remains only to establish that the symbolic growth rate of
the admissible grammar satisfies
\[
\lambda < 3
\]
The remainder of the paper is devoted to proving this spectral inequality.


% ============================================================
\section{The Beatty Barrier}
\label{sec:beatty-barrier}
% ============================================================

The remaining task is to bound the symbolic growth rate
\[
\lambda = \limsup_{m\to\infty} N(m)^{1/m}.
\]

The key observation is that the symbolic growth is not free.
Every admissible extension of a division word is governed by
the Beatty increment sequence associated with
\[
B(m)=\lfloor m\log_2 3\rfloor+1.
\]

Consequently, the admissible branching process is not arbitrary.
Every symbolic extension is constrained by the same Beatty sequence that
governs the dyadic contraction threshold. Thus the symbolic branching
process and the cylinder refinement process are controlled by a common
underlying arithmetic.

% ------------------------------------------------------------
\subsection{Beatty-Controlled Growth}
% ------------------------------------------------------------

The Beatty increment sequence
\[
\delta(m)=B(m+1)-B(m)\in\{1,2\}
\]
governs two parallel constructions:

\begin{enumerate}
\item the dyadic refinement of admissible cylinders; and
\item the doubling events that generate the symbolic counts
      \(N(m)\) in the horizontal-sum construction of
      Appendix~\ref{app:horizontal-sum-table}.
\end{enumerate}

The second statement follows from the construction of \OEIS{A186009},
whose recursive doubling rule is governed by the increment sequence
\OEIS{A022921} ~\cite{Winkler2017}, equal to the first differences of
\OEIS{A020914}.

Thus the same Beatty increment sequence governs both admissible cylinder
refinement and admissible symbolic branching. The two constructions are
therefore synchronized by a common Sturmian grammar.

Because the Beatty increment sequence is Sturmian, it admits canonical
return words of lengths \(2,5,12,\ldots\). Their construction from the
continued-fraction expansion of \(\log_2 3\) is summarized in
Appendix~\ref{app:sturmian}. For the present section we require only two
facts:

\begin{enumerate}
\item the canonical return words exist; and
\item each has a well-defined average dyadic increment.
\end{enumerate}

The next subsection shows that, once the increment sequence is fixed,
the induced horizontal-sum dynamics is represented by positive linear
operators, allowing the local symbolic growth contributed by each return
word to be estimated independently.

% ------------------------------------------------------------
\subsection{Linear Operator Formulation}
% ------------------------------------------------------------

Let $v_n = (a_1(n),a_2(n),\dots,a_{k_n}(n))$ denote the generating row at
level $n$, where the length $k_n$ depends on the increment history.
Let
\[
S_n = \sum_i a_i(n)
\]
be its total sum. Away from terminal duplication, the recurrence
\[
a_i(n)=a_i(n-1)+a_{i-1}(n)
\]
is linear with nonnegative coefficients. The terminal duplication is also
linear, so the complete evolution is represented by positive linear operators.

Representing the horizontal-sum dynamics by positive linear operators permits the
growth contributed by finite increment blocks to be analyzed through the corresponding
operator products.

A single increment acts as a positive linear endomorphism
$T_s^{(k)} : \mathbb{R}^{k} \to \mathbb{R}^{k}$, preserving row length, whereas
a double increment acts as a positive linear map
$T_d^{(k)} : \mathbb{R}^{k} \to \mathbb{R}^{k+1}$, increasing the row length by one
via terminal duplication.

The essential feature is positivity: every entry of \(T_s\) and \(T_d\) is nonnegative.
Consequently every admissible evolution is represented by a positive operator product.

Thus each increment symbol $\sigma_n\in\{s,d\}$ induces a positive linear map
\[
v_{n+1}=T_{\sigma_n}v_n,
\]
so every admissible division word corresponds to a finite product
\[
T_{\sigma_{n+k-1}}\cdots T_{\sigma_n}
\]
of positive linear operators.

\begin{lemma}[Controlled Growth of Horizontal Sums]
\label{lem:controlled-horizontal-growth}

Let a row have strictly positive entries with total sum $S$.
A single increment preserves the total row sum.

If a double increment is applied immediately afterward, producing a row in which
the terminal entry equals the sum of the previous row and all other entries remain
strictly positive, then the total sum of the resulting row satisfies
\[
S_{\mathrm{new}} < 3S.
\]
\end{lemma}

\begin{proof}
After the intermediate single-increment evolution, the terminal entry equals
\(S\), while the remaining interior entries are strictly positive and therefore
have total \(S_{\rm int}<S\).

The subsequent double increment duplicates the terminal entry, producing two
terminal contributions equal to \(S\). Hence

\[
S_{\rm new}
=
S_{\rm int}+2S,
\]
where \(S_{\rm int}<S\). Therefore
\[
S_{\rm new}<3S.
\]
\end{proof}

Since every admissible evolution is represented by positive operator products,
each primitive return block determines a finite operator product whose growth
factor depends only on that block. Consequently the symbolic contribution of
each primitive return word can be bounded locally.

Appendix~\ref{app:local-growth-details} computes these local growth factors for
every primitive return block. These provide uniform upper bounds for the symbolic
growth contributed by each admissible primitive block.


% ------------------------------------------------------------
\subsection{Beatty Block Drift}
% ------------------------------------------------------------

\begin{definition}[Dyadic Magnitude]
For a finite increment block $\Omega$, let
\[
\Delta(\Omega)
\]
denote the sum of its dyadic increments.
Equivalently, $\Delta(\Omega)$ is the total exponent of $2$
contributed by the block.
\end{definition}

\begin{lemma}[Beatty Drift Lemma]
\label{lem:beatty-drift-lemma}
Let \(\Omega\) be a finite block whose average increment satisfies
\[
\frac{\Delta(\Omega)}{|\Omega|} > \log_2 3
\]
Then the infinite periodic word
\[
\Omega^{\infty}
\]
is not admissible.
\end{lemma}
\begin{proof}
We define \(\Omega^k\) as the \(k\)-fold catentation of \(\Omega\).
This after \(k\) repetitions of \(\Omega\),
\[
\Delta(\Omega^k) = k \Delta(\Omega)
\]
while
\[
B(k|\Omega|) = k |\Omega| \log_2 3 + O(1)
\]
Since
\[
\Delta(\Omega) - |\Omega| \log_2 3 > 0,
\]
the difference
\[
\Delta(\Omega^k) - B(k|\Omega|)
\]

tend to \(+\infty\) as \(k \to \infty\). Since the excess drift grows
linearly while the Beatty discrepancy remains bounded, sufficiently long
prefixes violate the admissibility constraints. Therefore the infinite
periodic word \(\Omega^{\infty}\) is not admissible.
\end{proof}

\begin{remark}
By way of an explanation, \(\Omega\) posesses positive drift relative to
\(\log_2 3\). Anything with positive drift cannot persist forever because
the Beatty sequence never departs from its limiting slope by more than a
bounded amount whereas positive drift accumulates linearly without bound.
\end{remark}

\begin{lemma}[Extremal Admissibility Principle]
\label{lem:extremal-admissibility-principle}
Let \(\omega\) be an admissible infinite division word. Then
\[
\limsup_{n\to\infty} \frac{\Delta_\omega(n)}{n} \le \log_2 3
\]
Consequenty no admissible infinite word may contain an asymptotic density
of double increments exceeding the Beatty slope.
\end{lemma}

\begin{proof}
Suppose
\[
\limsup_{n\to\infty} \frac{\Delta_\omega(n)}{n} > \log_2 3.
\]

Then there exists
\(\varepsilon>0\)
and infinitely many prefixes satisfying
\[
\Delta_\omega(n) \ge (\log_2 3+\varepsilon)n.
\]

Hence these prefixes accumulate positive Beatty drift growing linearly
with \(n\), contradicting the Beatty Drift Lemma~\ref{lem:beatty-drift-lemma}.
\end{proof}


% ------------------------------------------------------------
\subsection{Canonical Return Words}
% ------------------------------------------------------------

The Beatty Drift Lemma~\ref{lem:beatty-drift-lemma} constrains only the
asymptotic average dyadic increment of admissible infinite words. To convert
this asymptotic restriction into a finite combinatorial statement, we use the
canonical return words of the Beatty Sturmian sequence.

The continued fraction expansion of \(\log_2 3\) yield the coefficents, \(a_k\) for
generating the Sturmian return words according to the recursive equation
\[
q_{k} = a_k q_{k-1} + q_{k-2}
\]

For the present argument, only the first two nontrivial canonical return words are
required, since every longer canonical return word is assembled recursively from them.

\begin{proposition}[Recursive Assembly of Canonical Return Words]
\label{prop:recursive-assembly}

Let \(C_2=(1,2)\) and
\(C_5=(1,2,1,2,2)\)
denote the first two nontrivial canonical Sturmian return words.
Every subsequent canonical return word is obtained recursively by
concatenating copies of \(C_2\) and \(C_5\).
\end{proposition}

\begin{proof}

The canonical Sturmian return words satisfy the standard continued-fraction
recursion
\[
C_k=C_{k-1}^{\,a_k}C_{k-2},
\]

where \(a_k\) is the corresponding continued-fraction coefficient of
\(\log_2 3\). The next composite word is
\[
C_{12}=C_2\,C_{5}^2
\]

Suppose inductively that both \(C_{k-1}\) and \(C_{k-2}\)
are concatenations of \(C_2\) and \(C_5\). Then
\[
C_k=C_{k-1}^{\,a_k}C_{k-2}
\]

is likewise a concatenation of the same primitive words.
The result follows by induction.
\end{proof}


% ------------------------------------------------------------
\subsection{The Beatty Barrier Theorem}
% ------------------------------------------------------------

The Beatty increment sequence satisfies
\[
a(n) = \lfloor (n+1) \log_2(3) \rfloor - \lfloor n \log_2 3 \rfloor
\]
and therefore
\[
\lim_{m \to \infty} \frac{1}{m} \sum_{k=1}^{m} a(k) = \log_2 3
\]
as recorded in \OEIS{A020914}. Among the primitive return words, the block
\[
C_5 = (1,2,1,2,2)
\]
has an avergage increment
\[
\frac{8}{5} > \log_2 3.
\]

Hence the Beatty Drift Lemma implies that the infinite periodic word
\[
(1,2,1,2,2)^\infty
\]
is not admissible. The obstruction is purely arithmetic: the excess average
increment accumulates linearly, while the Beatty discrepancy remains bounded.

\begin{theorem}[Beatty Barrier Theorem]
\label{thm:beatty-barrier}

No admissible infinite division word can contain the primitive block
\[
C_5 = (1,2,1,2,2)
\]
with asymptotic frequency one.  Equivalently, every admissible infinite
word must contain infinitely many longer return blocks whose average
dyadic increment is strictly smaller than $8/5$.
\end{theorem}

\begin{proof}
Suppose the asymptotic frequency is one. Then
\[
\frac{\Delta(n)}{n} \to \frac{8}{5}
\]
since
\[
\frac{8}{5} > \log_2 3.
\]
this contradicts the Extremal Admissibility Principle
(Lemma~\ref{lem:extremal-admissibility-principle}),
which asserts that every admissible infinite division word satisfies
\[
\limsup_{n\to\infty} \frac{\Delta_\omega(n)}{n} \le \log_2 3
\]

Therefore \(C_5\) cannot occur with asymptotic frequency one.
\end{proof}


% ============================================================
\section{Spectral Growth Bound}
\label{sec:spectral-growth-bound}
% ============================================================

The admissible grammar is constrained by the Beatty growth law \(B(n)\).
Any symbolic process whose average increment exceeds the Beatty slope must
eventually violate admissibility. The periodic 5-cycle is the minimal
non-trivial expansive process with this property.

Appendix~\ref{app:local-growth-details} computes the comparison growth
contributed by every primitive return block. Among these,
\[
C_5 = (1,2,1,2,2)
\]

has the largest per-symbol comparison growth.  Nevertheless,
\[
\Lambda_5 = \left( 3^2 \cdot
\frac{123}{33} \cdot
\frac{341}{131} \cdot
\frac{329}{123}
\right)^{1/5} < 3.
\]

\begin{lemma}[Extremal Block Dilution]
\label{lem:extremal-block-dilution}
Every occurrence of an expansive primitive contributes positive Beatty drift.
Since total drift must remain bounded, expansive primitives must be balanced
by sufficiently many contractive return words. Consequently their asymptotic
frequency is bounded away from one.
\end{lemma}

\begin{proof}
By the Extremal Admissibility Principle
(Lemma~\ref{lem:extremal-admissibility-principle}),
every admissible infinite division word has average dyadic increment at most
\(\log_2 3\).

The primitive block \(C_5\) has average increment \(8/5>\log_2 3\), so each
occurrence contributes positive Beatty drift relative to the admissible mean.

If the asymptotic frequency of expansive blocks approached one, the average
drift would exceed the Beatty slope, contradicting the Extremal Admissibility
Principle.

Hence expansive primitives must be balanced by return blocks having smaller
average increment. Consequently the asymptotic frequency of expansive blocks
is bounded away from one.
\end{proof}

\begin{theorem}[Spectral Growth Bound]
\label{thm:spectral-growth}
Every admissible symbolic trajectory satisfies
\[
\lambda < 3.
\]
\end{theorem}

\begin{proof}
The Beatty Barrier Theorem excludes infinite repetition of the unique
extremal primitive block.

Every admissible trajectory therefore contains infinitely many return
blocks with strictly smaller average growth.  By the Spectral Compensation
Lemma~\ref{lem:extremal-block-dilution} any instance of the 5-block one must
insert at least one compensating block.

Since every compensating block has comparison growth strictly smaller than
the extremal 5-block, the asymptotic exponential growth rate is strictly
below the growth obtained by perpetual repetition of \(C_5\).  Then
\[
\Lambda_5<3.
\]
Thus every admissible symbolic trajectory satisfies
\[
\lambda<3.
\]
\end{proof}

% ------------------------------------------------------------
\subsection{Interpretation}
% ------------------------------------------------------------

The Sturmian analysis identifies the primitive return words that may appear
in the Beatty grammar. The proof of the spectral bound depends only on the Beatty
drift restriction established in the previous section. The continued-fraction and
return-word analysis serves only to justify that the return cycles are canonical
primitive comparison words rather than ad hoc choices.

The Beatty-controlled affine grammar admits an extremal local growth mechanism
whose infinite repetition is structurally forbidden, and that forbidden mechanism
is already subcritical, namely \(\lambda < 3\).

Thus the measure of the residual forest collapses as \(n \to \infty\).


% ============================================================
\section{Conclusion}
% ============================================================

This paper reformulates the accelerated Collatz problem in terms of an admissible symbolic grammar whose
arithmetic realization is given by nested dyadic cylinders. The resulting cylinder forest provides an exact
symbolic representation of the odd integers together with a natural measure-theoretic filtration.

Within this framework the measure of the residual forest is controlled by two competing
exponential processes: symbolic proliferation and dyadic contraction. The measure-theoretic
analysis shows that these combine through a single spectral parameter $\lambda$. Consequently, proving
\[
\lambda < 3
\]
is sufficient to force the residual forest to have natural density zero.

Combining the spectral estimate with the residual density bounds shows that the
infinite residual set has natural density zero. Equivalently, the first-passage
layers exhaust the entire normalized dyadic measure.

The principal contribution is therefore not merely a new estimate, but a reformulation
of the density problem itself. Rather than studying individual Collatz trajectories
directly, the argument transfers the problem through three equivalent descriptions:
\[
\text{symbolic grammar} \longrightarrow \text{cylinder geometry} \longrightarrow \text{measure dynamics}.
\]
Within this framework, the density question becomes a problem of controlling the spectral
growth of an admissible Sturmian grammar.

This synthesis of symbolic dynamics, Diophantine approximation, and dyadic measure provides
a unified framework in which the density question for the accelerated Collatz map reduces
to a single explicit spectral inequality.


\appendix

% ============================================================
\newpage
\section{The Canonical Affine Map}
\label{app:canonical-c}
% ============================================================

Each division word \(\omega\) induces an affine map \(F_\omega(x)\) of length $m$
(see Section~\ref{sec:division-counts})
\[
F_\omega(x) = \frac{3^m x + c(\omega)}{2^{D(\omega)}}, \qquad D(\omega) = d_0 + \dots + d_{m-1},
\]
where \(c(\omega)\) is the affine constant canonically determined by the
division word \(\omega\).  The associated congruence constraint is naturally
considered modulo \(2^{D(\omega)}\).

Although the existence and uniqueness of \(c(\omega)\) follows from the affine representation
developed in Section~\ref{sec:division-counts}, the explicit computation is useful for concrete
examples and software implementations. The following procedure reconstructs the same affine
constant \(c(\omega)\) from any representative of the associated residue class.

\subsection{Canonical Choice}

For a fixed \(\omega\), any integer \(x\) in the corresponding residue class satisfies
\[
3^m x + c \equiv 0 \pmod{2^{D(\omega)}}.
\]
We remove ambiguity by selecting the unique representative
\[
0 \le c(\omega) < 2^{D(\omega)}.
\]
This choice depends only on \(\omega\) and is independent of the representative \(x\).

\subsection{Step-by-Step Computation}

Given \(\omega\) of length \(m\) and total division count \(D(\omega)\):

\begin{enumerate}
\item Choose any representative \(x\) of the residue class determined by \(\omega\).  
\item Compute \(y = 3^m x\).  
\item Let \(r = \lceil y / 2^{D(\omega)} \rceil\) be the smallest integer such that \(r \cdot 2^{D(\omega)} \ge y\).  
\item Define \(c(\omega) = r \cdot 2^{D(\omega)} - y\).  
\end{enumerate}

By this construction, \(0 \le c(\omega) < 2^{D(\omega)}\) ensuring that
\[
F_\omega(x) = \frac{3^m x + c(\omega)}{2^{D(\omega)}} \in \mathbb{Z} \quad \text{for all } x \text{ in the class.}
\]

\newpage
% ------------------------------------------------------------
\subsection{Example}
% ------------------------------------------------------------

Let \(\omega = (1,1,1,4)\) with \(m=4\) and \(D(\omega) = \sum_{i=0}^{3} d_i = 1 + 1 + 1 + 4 = 7\):

\begin{enumerate}
\item Choose \(x = 15\). Compute \(y = 3^4 \cdot 15 = 1215\).  
\item \(2^7 = 128\), so \(r = \lceil 1215/128 \rceil = 10\).  
\item \(c(\omega) = 10 \cdot 128 - 1215 = 65\).  
\end{enumerate}

Hence the canonical map is
\[
F_\omega(x) = \frac{81x + 65}{128}.
\]

Checking another representative \(x = 143\):
\[
F_\omega(143) = \frac{81 \cdot 143 + 65}{128} = 91,
\]
confirming that the same \(c(\omega)\) works for the entire residue class.
The admissibility condition becomes
\[
81x+65\equiv0\pmod{128}.
\]

Solving yields the unique residue class
\[
x\equiv15\pmod{128}.
\]

Hence the cylinder is
\[
x(k)=15+128k.
\]

Applying the affine map,
\[
F_\omega(15+128k)
=
81k+10.
\]

Thus the symbolic cylinder generated by \(\omega\) is mapped onto the
arithmetic progression
\[
10,\,91,\,172,\,253,\dots
\]
with common difference \(81=3^4\).

% ------------------------------------------------------------
\subsection*{Remark}
% ------------------------------------------------------------

This construction is purely illustrative and is not required for the density
or growth arguments in the main text. It demonstrates concretely how the affine
map is fully determined by the division word \(\omega\).

By the Cylinder Representation Theorem (Section~\ref{sec:cylinder-geometry}), every
realization of \(\omega\) belongs to a single residue class modulo \(2^{D(\omega)}\).
Consequently, the value of \(c(\omega)\) computed above is independent of the chosen
representative.

\begin{remark}[Cylinder Rigidity]
The division word $\omega$ uniquely determines its canonical affine
constant $c(\omega)$, and therefore uniquely determines the affine map
\[
F_\omega(x) = \frac{3^m x+c(\omega)}{2^{D(\omega)}}.
\]
Since admissible realizations of $\omega$ occupy a single residue class
modulo $2^{D(\omega)}$, the symbolic data encoded by $\omega$
simultaneously determines both the affine dynamics and the associated
cylinder.
\end{remark}


% ------------------------------------------------------------
\subsection{Illustrative Derivation of the Affine Constant}
% ------------------------------------------------------------

For readers wishing to see the affine constant arise directly from
iterated substitution, consider a general division word
\[
\omega=(a,b,c,d)
\]
of length \(4\).

The accelerated relations are
\[
x_1=\frac{3x_0+1}{2^a},
\qquad
x_2=\frac{3x_1+1}{2^b},
\qquad
x_3=\frac{3x_2+1}{2^c},
\qquad
x_4=\frac{3x_3+1}{2^d}.
\]

Eliminating intermediate variables recursively yields
\[
2^{a+b+c+d}x_4
=
3^4x_0
+
3^3
+
3^2 2^a
+
3\,2^{a+b}
+
2^{a+b+c}.
\]

Hence
\[
F_\omega(x)
=
\frac{
3^4x
+
\left(
3^3
+
3^2 2^a
+
3\,2^{a+b}
+
2^{a+b+c}
\right)
}{
2^{a+b+c+d}
},
\]
so that
\[
c(\omega)
=
3^3
+
3^2 2^a
+
3\,2^{a+b}
+
2^{a+b+c}.
\]

This illustrates concretely how the affine constant accumulates the
propagated contributions of successive ``+1'' terms throughout the
accelerated iteration.

% ============================================================
\newpage
\section{Continued Fractions, Tower Cycles, and the Extremal Envelope}
\label{app:continued-fractions}
% ============================================================

This appendix records the Diophantine structure underlying the admissible
grammar and explains the origin of the distinguished cycle lengths
\[
1,\ 1,\ 2,\ 5,\ 12,\ 41,\ 53,\ 306,\ 665,\ 15601,\ 31867,\ 79335,\ldots
\]
that appear throughout the paper.

% ------------------------------------------------------------
\subsection{Continued-Fraction Origin of the Tower Cycles}
% ------------------------------------------------------------

Let

\[
\alpha=\log_2 3.
\]

Its continued-fraction expansion begins

\[
\alpha=[1;1,1,2,2,3,1,5,2,23,2,2,\dots].
\]

Writing

\[
\frac{p_k}{q_k}
\]

for the convergents of $\alpha$, the first denominator sequence is

\[
1,\ 1,\ 2,\ 5,\ 12,\ 41,\ 53,\ 306,\ 665,\ 15601,\ldots
\]

which coincides exactly with the tower lengths observed in the admissible grammar.
The convergent denominators satisfy the standard recurrence

\[
q_k=a_k q_{k-1}+q_{k-2},
\]

where $a_k$ is the corresponding continued-fraction coefficient.
Consequently every tower cycle is assembled canonically from the two
preceding tower cycles.

Table~\ref{tab:tower-cycles} records the first instances.

\begin{table}[ht]
\centering
\caption{Continued-fraction generation of the tower cycles.}
\label{tab:tower-cycles}
\begin{tabular}{r l r r l l}
\hline
$q_k$ & Cycle & $k$ & $a_k$ & Cycle Composition & Sequence \\
\hline
1     & $C_{    0}$ &  0 & -- & Single                          & $1$ \\
1     & $C_{    1}$ &  1 &  1 & Double                          & $2$ \\
2     & $C_{    2}$ &  2 &  1 & $C_{    0} + 1 \cdot C_{    1}$ & $1,2$ \\
5     & $C_{    5}$ &  3 &  2 & $2 \cdot C_{    2} + C_{    1}$ & $1,2,1,2,2$ \\
12    & $C_{   12}$ &  4 &  2 & $C_{    2} + 2 \cdot C_{    5}$ & $1,2,1,2,1,2,2,1,2,1,2,2$ \\
41    & $C_{   41}$ &  5 &  3 & $3 \cdot C_{   12} + C_{    5}$ & $1,2,1,2,1,2,2,\ldots,1,2,1,2,2$  \\
53    & $C_{   53}$ &  6 &  1 & $C_{   12} + 1 \cdot C_{   41}$ & $1,2,1,2,1,2,2,\ldots,1,2,1,2,2$  \\
306   & $C_{  306}$ &  7 &  5 & $5 \cdot C_{   53} + C_{   41}$ & $1,2,1,2,1,2,2,\ldots,1,2,1,2,2$  \\
665   & $C_{  665}$ &  8 &  2 & $C_{   53} + 2 \cdot C_{  306}$ & $1,2,1,2,1,2,2,\ldots,1,2,1,2,2$  \\
15601 & $C_{15601}$ &  9 & 23 & $23 \cdot C_{665} + C_{   306}$ & $1,2,1,2,1,2,2,\ldots,1,2,1,2,2$  \\
31867 & $C_{31867}$ & 10 &  2 & $C_{  665} + 2 \cdot C_{15601}$ & $1,2,1,2,1,2,2,\ldots,1,2,1,2,2$  \\
79335 & $C_{79335}$ & 11 &  2 & $2 \cdot C_{31867} + C_{15601}$ & $1,2,1,2,1,2,2,\ldots,1,2,1,2,2$  \\
\hline
\end{tabular}
\end{table}

Thus the tower hierarchy is not an empirical feature of the grammar.
It is forced by the Diophantine approximation structure of $\log_2 3$.

% ------------------------------------------------------------
\subsection{Alternating Expansion and Contraction}
% ------------------------------------------------------------

The convergents alternately overestimate and underestimate $\alpha$.

Define
\[
\varepsilon_k = q_k\alpha-p_k.
\]

Then
\[
\varepsilon_k \to 0
\]
and the signs alternate.

Since
\[
\alpha=\log_2 3,
\]
it follows that
\[
\varepsilon_k = q_k\log_2 3-p_k = -\log_2\!\left(\frac{2^{p_k}}{3^{q_k}}\right).
\]

Equivalently,
\[
\frac{2^{p_k}}{3^{q_k}} = 2^{-\varepsilon_k}.
\]

Hence

\[
\left|
\log\!\left(
\frac{2^{p_k}}{3^{q_k}}
\right)
\right|
=
(\log 2)\,|\varepsilon_k|
\to0.
\]

Therefore

\[
\frac{2^{p_k}}{3^{q_k}}
\to1.
\]

The convergents alternately overestimate and underestimate
$\log_2 3$, so the ratios alternately lie above and below unity while
approaching exact multiplicative balance.

\begin{table}[ht]
\centering
\caption{Convergents of $\log_2 3$ and resulting
dyadic--ternary synchronization ratios.}
\label{tab:tower-ratios}
\begin{tabular}{r r c c}
\hline
$q_k$ & $p_k$ &
$\varepsilon_k=q_k\log_2(3)-p_k$ &
$\displaystyle 2^{p_k} / 3^{q_k}$ \\
\hline
 1 &  2 & $-0.41594\ldots$ & $1.33333\ldots$ \\
 2 &  3 & $+0.16993\ldots$ & $0.88889\ldots$ \\
 5 &  8 & $-0.07519\ldots$ & $1.05350\ldots$ \\
12 & 19 & $+0.01955\ldots$ & $0.98654\ldots$ \\
41 & 65 & $-0.01654\ldots$ & $1.01153\ldots$ \\
53 & 84 & $+0.00301\ldots$ & $0.99791\ldots$ \\
\hline
\end{tabular}
\end{table}

% ------------------------------------------------------------
\subsection{Primitive Return Words and the Failure of the Double Process}
% ------------------------------------------------------------

The Sturmian word generated by
\[
d(n)=\lfloor n\alpha\rfloor+1
\]
admits a return-word decomposition.  At the first nontrivial level, the primitive
return words have lengths $2$, $5$, and $12$. These correspond respectively to:

\[
C_2=(1,2).
\]
\[
C_5=(1,2,1,2,2).
\]
\[
C_{12}=(1,2,1,2,1,2,2,1,2,1,2,2).
\]

The 2-cycle is locally contractive, while the 5-cycle is locally expansive. A 12-cycle
is formed by concatenating the 2-cycle with a pair of consecutive 5-cycles, but the 
12-cycle is overall contractive.

One might ask why the extremal comparison process is based on the
5-cycle rather than the shorter expansive 1-cycle corresponding to a
double increment.

Indeed, the shortest expansive block is

\[
C_1=(2),
\]

for which

\[
\frac{2^2}{3} = \frac43 > 1.
\]

Repeated doubles generate the row sums

\[
1,2,5,14,42,132,429,\ldots
\]

which are the Catalan numbers.
Since

\[
C_n
\sim
\frac{4^n}{n^{3/2}\sqrt{\pi}}
\]

implies

\[
\lim_{n\to\infty}
\frac{C_{n+1}}{C_n}
=
4.
\]

Thus an infinite concatenation of doubles produces exponential growth
rate \(4\), which exceeds the critical value \(3\).

The pure-double process is therefore unsuitable as a comparison
envelope: it overestimates growth by too large a margin and cannot
yield the desired spectral bound.

% ------------------------------------------------------------
\subsection{The Extremal Envelope}
% ------------------------------------------------------------

The main text proves that every admissible Sturmian word decomposes uniquely
into primitive return words. Since the normalized cocycle contribution of \(C_5\)
exceeds that of every other primitive return word, replacing each return word by \(C_5\)
can only increase the average cocycle drift. Although the periodic comparison word
\(C_5^\infty\) is itself forbidden by the Sturmian grammar, it provides a universal
comparison envelope. The explicit evaluation of this envelope is carried out in 
Appendix~\ref{app:local-growth-details}.

% ============================================================
\newpage
\section{Extremal Local Growth Calculations}
\label{app:local-growth-details}
% ============================================================

This appendix provides the explicit finite computations underlying the local
growth bounds used in the density argument. Recall that the sequence \OEIS{A020914}
has first differences given by \OEIS{A022921}, consisting of repeated blocks of
\emph{single} steps ($s$) and \emph{double} steps ($d$).

% ------------------------------------------------------------
\subsection{Illustrative Extremal Computations}
% ------------------------------------------------------------

The generating rows evolve by a horizontal-sum rule:
\begin{equation}
a_i(n) = a_i(n-1) + a_{i-1}(n), \quad i,n \in \mathbb{N},
\end{equation}

so that sequences of single steps exhibit ordinary Pascal--type growth.
Double steps differ only in that the terminal entry of the row is duplicated,
causing the final term to appear twice in the generating tuple while preserving
monotonicity of the preceding entries.

\begin{lemma}[Pascal-Type Horizontal Growth]
\label{lem:pascal-growth}

Let a row generated by the horizontal-sum construction have strictly positive entries
with total sum $S$.

After one horizontal-sum step without duplication, the resulting row has total

\[
S_{\mathrm{new}}
=
S+S_{\mathrm{int}},
\]

where $S_{\mathrm{int}}$ denotes the sum of the nonterminal entries of the
previous row.

Consequently

\[
S_{\mathrm{new}} < 2S.
\]

\end{lemma}

\begin{proof}

The terminal entry of the new row equals the total sum $S$ of the previous row.
Every nonterminal entry contributes exactly once to the new row through the
horizontal-sum recurrence. Hence

\[
S_{\mathrm{new}}
=
S+S_{\mathrm{int}}.
\]

Since all entries are strictly positive,

\[
S_{\mathrm{int}}<S,
\]

and therefore

\[
S_{\mathrm{new}}<2S.
\]

\end{proof}

\begin{corollary}[Controlled Growth Under Duplication]
\label{lem:controlled-growth-under-duplication}

A subsequent duplication contributes at most one additional copy of the
terminal entry, whose value is $S$.

Hence

\[
S_{\mathrm{new}}
<
2S+S
=
3S.
\]

\end{corollary}

These examples are chosen so that each step saturates the inequalities of
Lemma~\ref{lem:pascal-growth}, thereby realizing extremal growth. Equalities
shown below should be interpreted as extremal upper-bound realizations rather than
identities valid for all rows.

The following row-by-row computations realize extremal growth for the admissible
local words and together determine the sharpest uniform bounds used in
Proposition ~\ref{prop:extremal-comparison-word}.  Each example illustrates how a local pattern
can saturate its maximal growth factor under the horizontal-sum dynamics.  Let $S$
denote the total sum of entries in the generating row prior to a given local pattern.

\paragraph{Example 1.}
\[
\begin{array}{l}
\textit{Single } a(10) = \underbrace{173}_{\text{sum} := S}:
\underbrace{1,7,25,55,85}_{S}
\end{array}
\]

\[
\begin{array}{l}
\textit{Double } a(11) = \underbrace{476}_{\text{sum} = 3S}:
\underbrace{1,8,33,88}_{S}
\underbrace{173}_{S},
\underbrace{173}_{S}
\end{array}
\]

\[
\begin{array}{l}
\textit{Single } a(12) = \underbrace{961}_{\text{sum} = 9S}:
\underbrace{1,9,42,130}_{2S}
\underbrace{303}_{3S},
\underbrace{476}_{4S}
\end{array}
\]

\[
\begin{array}{l}
\textit{Double } a(13) = \underbrace{2652}_{\text{sum} = 33S}:
\underbrace{1,10,52,182}_{4S},
\underbrace{485}_{7S},
\underbrace{961}_{11S},
\underbrace{961}_{11S}
\end{array}
\]

\[
\begin{array}{l}
\textit{Double } a(14) = \underbrace{8045}_{\text{sum} = 123S}:
\underbrace{1,11,63,245}_{8S},
\underbrace{730}_{15S},
\underbrace{1691}_{26S},
\underbrace{2652}_{37S},
\underbrace{2652}_{37S}
\end{array}
\]

\[
\begin{array}{l}
\textit{Single } a(15) = \underbrace{17637}_{\text{sum} = 329S}:
\underbrace{1,12,75,320}_{16S},
\underbrace{1050}_{31S},
\underbrace{2741}_{57S},
\underbrace{5393}_{94S},
\underbrace{8045}_{131S}
\end{array}
\]

The resulting maximal local ratios are
\[
\begin{array}{c|c}
\text{Pattern} & \text{Max Ratio} \\
\hline
s,d & 3/1 \\
d,d & 123/33 \\
s,d,s & 9/3 \\
d,d,s & 329/123
\end{array}
\]

\paragraph{Example 2.}
\[
\begin{array}{l}
\textit{Single } a(12) = \underbrace{961}_{\text{sum} := S}:
\underbrace{1,9,42,130,303,476}_{S}
\end{array}
\]

\[
\begin{array}{l}
\textit{Double } a(13) = \underbrace{2652}_{\text{sum} = 3S}:
\underbrace{1,10,52,182,485}_{S},
\underbrace{961}_{S},
\underbrace{961}_{S}
\end{array}
\]

\[
\begin{array}{l}
\textit{Double } a(14) = \underbrace{8045}_{\text{sum} = 13S}:
\underbrace{1,11,63,245,730}_{2S},
\underbrace{1691}_{3S},
\underbrace{2652}_{4S},
\underbrace{2652}_{4S}
\end{array}
\]

\[
\begin{array}{l}
\textit{Single } a(15) = \underbrace{17637}_{\text{sum} = 37S}:
\underbrace{1,12,75,320,1050}_{4S},
\underbrace{2741}_{7S},
\underbrace{5393}_{11S},
\underbrace{8045}_{15S}
\end{array}
\]

\[
\begin{array}{l}
\textit{Double } a(16) = \underbrace{51033}_{\text{sum} = 131S}:
\underbrace{1,13,88,408,1458}_{8S},
\underbrace{4199}_{15S},
\underbrace{9592}_{26S},
\underbrace{17637}_{41S},
\underbrace{17637}_{41S}
\end{array}
\]

\[
\begin{array}{l}
\textit{Single } a(17) = \underbrace{108950}_{\text{sum} = 341S}:
\underbrace{1,14,102,510,1968}_{16S},
\underbrace{6167}_{31S},
\underbrace{15759}_{57S},
\underbrace{33396}_{98S},
\underbrace{51033}_{139S}
\end{array}
\]

The corresponding maximal local ratios are
\[
\begin{array}{c|c}
\text{Pattern} & \text{Max Ratio} \\
\hline
s,d & 3/1 \\
d,d & 13/3 \\
s,d,s & 341/131 \\
d,d,s & 37/13
\end{array}
\]

% ------------------------------------------------------------
\subsection*{Improved Extremal Alignment}
% ------------------------------------------------------------

Because the total multiplicative growth over a cycle is the product of the local
growth factors, the geometric mean over one full cycle is maximized by selecting
the maximal admissible factor at each local pattern occurrence. Consequently,
once a uniform upper bound has been established for each admissible local pattern,
the extremal global growth rate over a complete cycle is obtained by multiplying
these maximal local factors according to their frequencies and taking the
corresponding geometric mean.

\emph{Equivalently, because total growth over any admissible word is multiplicative in
its local factors, the asymptotic per-symbol growth rate is the geometric mean of these
factors;  hence the maximal rate over the entire language generated by the grammar is
achieved by the word that maximizes this mean subject to the grammar constraints.}

Each example yields a valid (not necessarily sharp) upper bound on the maximal growth
factor associated with each admissible local pattern.  Since these bounds arise from
the same uniform recurrence and apply globally to all admissible configurations, the
true maximal growth factor for a given pattern is bounded above by the smallest
recorded bound among all such valid upper bounds.

Because the recurrence is monotone with respect to the placement of delayed doubling,
no other admissible configuration can exceed these extremal realizations.
Accordingly, combining the sharpest bounds obtained across different initiating
rows yields the tightest global envelope on admissible local growth.

\[
\begin{array}{c|c|c}
\text{Pattern} & \text{Max Ratio} & \text{Associated coefficient count} \\
\hline
s,d & 3 & 2 \\
d,d & 123/33 & 1 \\
s,d,s & 341/131 & 1 \\
d,d,s & 329/123 & 1
\end{array}
\]

Consequently, the maximal geometric growth rate of a $5$--cycle satisfies
\begin{equation}
\Bigl(
3^2 \cdot (123/33) \cdot (341/131) \cdot (329/123)
\Bigr)^{1/5} 
\approx 2.976 < 3.
\label{eq:rho-5}
\end{equation}

\paragraph{Conclusion.}
Every admissible growth pattern of \OEIS{A186009} is constrained by the grammar
of \OEIS{A022921}. Infinite concatenations of 5-cycles, though forbidden, saturate
local growth constraints and provide a strict upper envelope. Equation~\eqref{eq:rho-5}
confirms that even this extremal envelope has geometric growth rate $< 3$, ensuring
subdominance relative to the denominator $2^{B(i)}$.


% ============================================================
\section{Extremal Growth in the Sturmian Return-Word System}
\label{app:sturmian}
% ============================================================

% ------------------------------------------------------------
\subsection{Dyadic Threshold Sequence Sturmian Structure}
% ------------------------------------------------------------

\begin{theorem}[Sturmian realization of the threshold grammar]
Let
\[
B(n)=\lfloor n\log_2 3\rfloor+1.
\]
Then the increment sequence
\[
\delta(n)=B(n+1)-B(n)
\]

is a Sturmian word over the alphabet \(\{1,2\}\).
Consequently the admissible grammar inherits the combinatorial
structure of an irrational rotation.
\end{theorem}

\begin{proof}
This is the classical characterization of Beatty increment sequences
associated with irrational slopes.
\end{proof}

Let
\[
\alpha := \log_2 3,
\qquad
d(n) := \lfloor n\alpha \rfloor + 1.
\]
The increment sequence
\[
\delta(n) := d(n+1) - d(n) \in \{1,2\}
\]
is a Sturmian word of slope $\alpha - 1$ and intercept $0$.
Equivalently, $\delta$ defines a symbolic dynamical system
\[
X_\alpha \subset \{1,2\}^{\mathbb{N}}
\]
generated by an irrational rotation. All admissible growth structures
considered in this paper are functions of this Sturmian coding.

% ------------------------------------------------------------
\subsection{Return Words and Primitive Cycle Decomposition}
% ------------------------------------------------------------

\begin{definition}[Return Words]
A finite word $C$ is a return word of the Sturmian system $X_\alpha$ if it is
the minimal block between successive visits to a fixed cylinder set induced by
$\delta$. We denote by $\mathcal{R}_\alpha$ the set of return words.
\end{definition}

\begin{theorem}[Return Words via Continued Fractions]
Let $\alpha = \log_2 3$ and let $\delta(n)$ be the Sturmian coding induced
by rotation $R_\alpha$.

Let $q_k$ denote the denominators of the continued-fraction convergents of $\alpha$.
For the prefix defining the extremal admissible growth process, the set of return
words is exactly:
\[
\mathcal{R}_\alpha = \{C_{q_k}, C_{q_{k+1}}\}
\]
for some index $k$ determined by the chosen prefix.  In particular, for the admissible
threshold structure of $d(n)=\lfloor n\alpha \rfloor +1$, the relevant convergent pair is:
\[
(q_k, q_{k+1}) = (5,12).
\]
\end{theorem}

\begin{proof}
The sequence $\delta(n)$ is a Sturmian word of slope $\alpha$.
Return words to any prefix correspond to first-return maps of the rotation
$R_\alpha$ to a subinterval $I$.

The convergent denominators begin (see Appendix \ref{app:continued-fractions} for details):
% The convergent denominators begin:
\[
1,\ 1,\ 2,\ 5,\ 12,\ 41,\ 53,\ 306,\ 665,\ 15601,\ 31867,\ 79335,\ldots
\]

For the prefix relevant to the admissible grammar, the associated return times are the
consecutive denominators 5 and 12. Hence the return-word set is $\{C_5, C_{12}\}$.
\end{proof}

Although the shortest return word is the single double increment $C_1=(2)$, repeated doubles generate the
Catalan sequence $1,2,5,14,42,132,\ldots$ whose asymptotic exponential growth rate is 4. Since this
exceeds the critical threshold 3, the pure-double process cannot serve as a useful comparison envelope.

\begin{proposition}[Extremal Primitive Comparison Word]
\label{prop:extremal-comparison-word}
The return word \(C_5\) is the shortest return word whose induced
comparison growth is strictly less than \(3\).
\end{proposition}

\begin{remark}[Diophantine Origin of the Grammar]
\label{rem:diophantine-grammar}

The admissible grammar is governed not by arbitrary combinatorial rules but by the Diophantine
approximation properties of \(\log_2 3\).  The convergent denominators
\[
1,\ 1,\ 2,\ 5,\ 12,\ 41,\ 53,\ 306,\ 665,\ 15601,\ 31867,\ 79335,\ldots
\]
arise from the continued-fraction approximants of
\[
\alpha=\log_2 3,
\]
and consequently, the primitive return words governing admissible refinement are
determined by best rational approximations to the Beatty slope rather than by
ad hoc combinatorial choices.
\[
\frac{p_k}{q_k}\approx \log_2 3.
\]

Equivalently,
\[
2^{p_k}\approx 3^{q_k},
\]
so the associated multiplicative cocycle satisfies
\[
\frac{2^{p_k}}{3^{q_k}}\to1.
\]

In particular, the appearance of the 5-, 12-, 41-, and 53-cycles
is not accidental; these cycles are induced by the Diophantine
approximation hierarchy of $\log_2 3$.
\end{remark}

\begin{lemma}[Induced Return-Word Hierarchy]
\label{lem:induced-hierarchy}

Every admissible word in the Sturmian language generated by
\[
\delta(n) = \lfloor (n+1)\alpha\rfloor-\lfloor n\alpha\rfloor
\]
admits a decomposition into return words generated
recursively from the primitive pair $(C_5, C_{12})$.
\end{lemma}

\noindent

% ------------------------------------------------------------
\subsection{Multiplicative Cocycle Structure}
% ------------------------------------------------------------

Define the multiplicative growth cocycle
\[
G(w) := \frac{2^{\#s(w) + 2\#d(w)}}{3^{|w|}},
\]
where $\#s(w)$ and $\#d(w)$ are the number of single and double increments in $w$.

\begin{lemma}[Cocycle Additivity]
For concatenation of admissible return words,
\[
G(C_{i_1} \cdots C_{i_n}) = \prod_{k=1}^n G(C_{i_k}).
\]
Equivalently,
\[
\log G(w) = \sum_{k} \log G(C_{i_k}).
\]
\end{lemma}
Thus the global exponential growth rate reduces to weighted averages
of return-word contributions.

% ------------------------------------------------------------
\subsection{Local Evaluation of Return Words}
% ------------------------------------------------------------

The local comparison constants computed in Appendix \ref{app:local-growth-details}
yield the following per-symbol cocycle contributions. The two primitive return words satisfy:

\[
\frac{1}{5}\log G(C_5)
=
\log\!\left(
3^{2}
\cdot \frac{123}{33}
\cdot \frac{341}{131}
\cdot \frac{329}{123}
\right)^{1/5},
\]

\[
\frac{1}{12}\log G(C_{12})
<
\frac{1}{5}\log G(C_5).
\]

Hence $C_5$ is the unique extremal return word with maximal per-symbol drift.

\begin{proposition}[Extremal Cocycle Principle]
For any admissible infinite word $\delta \in X_\alpha$,
\[
\limsup_{n\to\infty}
\frac{1}{n} \log G(\delta_{[1,n]})
\le
\frac{1}{5}\log G(C_5).
\]
\end{proposition}

\begin{proof}
By the unique decomposition into return words and cocycle additivity,
the growth rate is a convex combination of the per-symbol contributions
of $C_5$ and $C_{12}$.

Since
\[
\frac{1}{5}\log G(C_5)
>
\frac{1}{12}\log G(C_{12}),
\]
the maximal comparison average is obtained by replacing every return word with $C_5$,
corresponding formally to the periodic comparison word $C_5^\infty$.
\end{proof}

% ------------------------------------------------------------
\subsection{Global Growth Bound}
% ------------------------------------------------------------

\begin{theorem}[Spectral Growth Bound]
\label{thm:spectral-growth-bound}
Let $\lambda$ denote the symbolic growth rate introduced in Section
\ref{sec:symbolic-growth}.

Then
\[
\lambda \le \exp\!\left(\frac{1}{5}\log G(C_5)\right) < 3.
\]
\end{theorem}

\begin{proof}
% From the extremal return-word principle, the maximal growth is achieved by
% periodic concatenation of $C_5$. Evaluating the local constants yields
By Proposition~\ref{prop:extremal-comparison-word}, every admissible infinite
word has average cocycle growth bounded above by that of \(C_5^\infty\).  Thus
\[
\lambda \le \exp\!\left(\frac{1}{5}\log G(C_5)\right).
\]
Via explicit evaluation of the local cocycle contributions (Appendix \ref{app:local-growth-details}),
\[
\left(
3^{2}
\cdot \frac{123}{33}
\cdot \frac{341}{131}
\cdot \frac{329}{123}
\right)^{1/5}
< 3.
\]
\end{proof}
Theorem~\ref{thm:spectral-growth-bound} closes the logical chain developed throughout
the paper. The Beatty threshold sequence induces the admissible grammar; the grammar
determines the symbolic growth rate \(\lambda\); the Sturmian return-word analysis
establishes the strict bound \(\lambda < 3\); and Theorem~\ref{thm:residual-density-collapse}
then forces the residual forest to have zero natural density. Thus the symbolic, arithmetic,
and measure-theoretic components of the argument fit together in a single chain of implications.


\newpage
\section{Horizontal-Sum Table for A186009}
\label{app:horizontal-sum-table}

For reference, a portion of \OEIS{A186009} generated by the horizontal-sum construction is
given below as per an algorithm derived from Winkler \cite{Winkler2017}. This table is included
for illustration of the horizontal-sum mechanism; none of the density or growth bounds in the
main text depend on these specific numeric values.

\begin{table}[htbp]
\centering
\caption{A186009 Horizontal Sum Formulation}
\label{tab:A186009}
\vspace{4pt}
\begin{tabular}{@{}rrrrrrrrrrr@{}}
\toprule
$n$ & $a_1(n)$ & $a_2(n)$ & $a_3(n)$ & $a_4(n)$ & $a_5(n)$ &
$a_6(n)$ & $a_7(n)$ & $a_8(n)$ & $a_9(n)$ & $\sum a_i(n)$ \\
\midrule
1  & 1 &   &   &   &   &   &   &   &   & 1 \\
2  & 1 &   &   &   &   &   &   &   &   & 1 \\
3  & 1 &   &   &   &   &   &   &   &   & 1 \\
4  & 1 & 1 &   &   &   &   &   &   &   & 2 \\
5  & 1 & 2 &   &   &   &   &   &   &   & 3 \\
6  & 1 & 3 & 3 &   &   &   &   &   &   & 7 \\
7  & 1 & 4 & 7 &   &   &   &   &   &   & 12 \\
8  & 1 & 5 & 12 & 12 &   &   &   &   &   & 30 \\
9  & 1 & 6 & 18 & 30 & 30 &   &   &   &   & 85 \\
10 & 1 & 7 & 25 & 55 & 85 &   &   &   &   & 173 \\
11 & 1 & 8 & 33 & 88 & 173 & 173 &   &   &   & 476 \\
12 & 1 & 9 & 42 & 130 & 303 & 476 &   &   &   & 961 \\
13 & 1 & 10 & 52 & 182 & 485 & 961 & 961 &   &   & 2652 \\
14 & 1 & 11 & 63 & 245 & 730 & 1691 & 2652 & 2652 &   & 8045 \\
15 & 1 & 12 & 75 & 320 & 1050 & 2741 & 5393 & 8045 &   & 17637 \\
16 & 1 & 13 & 88 & 408 & 1458 & 4199 & 9592 & 17637 & 17637 & 51033 \\
17 & 1 & 14 & 102 & 510 & 1968 & 6167 & 15759 & 33396 & 51033 & 108950 \\
\bottomrule
\end{tabular}
\end{table}

There is a sequence closely related to \OEIS{A186009} called \OEIS{A100982}.  \OEIS{A186009}
prepends the number 1 to the \OEIS{A100982} sequence.  Along with the listing of initial
terms for \OEIS{A100982} comes an algorithm which explains how to generate the entire sequence
starting with an initial value of 1.  As a result of its close similarity, this method can
also be leveraged to generate all of the terms in \OEIS{A186009}. 

Theorem 2 \cite{Winkler2017} referenced in \OEIS{A100982} asserts that the elements in the
sequence can be derived in a Pascal–triangle–type manner, beginning with an initial condition:
\[
a(1) = a_1(1) = 1
\]

The algorithm for $A100982(n)$ states that the $k^{th}$ term of the $n^{th}$ series component
can be computed as follows:
% Recursive algorithm for generating new terms for A100982 based on previously generated terms
\[
a_{k+1}(n) = a_k(n) + a_k(n-1); k \in \setNo, n \in \setN
\]

The algorithm also stipulates that the last term in the sequence is repeated whenever
$A022921(n-2)=2$ (which can be derived from \OEIS{A020914} by taking the first
differences between its terms).  In doing so, this imposes the implicit constraint that this rule
applies only when $n \ge 2$.  Keeping in mind that repeated terms in \OEIS{A022921}
occur only once $n \ge 2$, the recurrence is extended to $n=1$ by treating the duplication
condition as inactive at the boundary, producing an augmented vertical-sum expansion.

In rows without terminal duplication, this recurrence is exactly Pascal’s binomial rule; duplication acts
only at the boundary and therefore preserves the monotone, binomial-type growth of interior entries.
This boundary-only modification is precisely why the extremal growth analysis in
Appendix~\ref{app:local-growth-details} reduces to \emph{local pattern frequencies}
rather than global row shape.

We reformulate this construction such that the first non-zero term in
each row is always at index 1 and thus the formula henceforth (producing the same expansion
components up to a reindexing of coordinates) is modified as follows.

Preserving $a(1) = a_1(1) = a_1(n) = 1 \text{ for all } n\in\setN$, the $i^{th}$ term of the
$n^{th}$ \OEIS{A100982} element is:
% Modified recursive algorithm for generating new terms for A100982 based on previously generated terms
\[
a_{i}(n) = a_i(n-1) + a_{i-1}(n); i,n \in \setN
\]

This reindexing does not change any row sums or terminal-duplication positions; it only relocates coordinates.
Using this formulation the initial term for all values of $n \in \setN$ is always 1 rather than
shifting out for each value of $n$.  The other significant difference is that in this formulation
the values for $a_i(n)$ appear horizontally as opposed to vertically.  The doubling rule for the
terminal entry remains unchanged in substance – namely that it occurs when $A022921(n-2) = 2$ -
which is precisely the single/double grammar governing admissible growth in the preceding appendices,
and hence the local multiplicative structure used in the extremal bounds.

Because \OEIS{A186009} is obtained from \OEIS{A100982} by prepending an initial row, all row indices
shift by one.  Consequently the terminal-duplication condition expressed in the original
formulation as $A022921(n-2)=2$ now becomes $A022921(n-3)=2$ in the present indexing.


% ============================================================
\clearpage
\bibliographystyle{unsrtnat}
\bibliography{src/latex/references}

\end{document}